\documentclass[11pt]{article}

\usepackage{amssymb,amsmath,amsfonts,amsthm}
\usepackage{latexsym}
\usepackage{graphics}
\usepackage{indentfirst}
\usepackage{lmodern}
\usepackage{hyperref}
\usepackage{mathrsfs}
\usepackage{array}
\usepackage{mathtools}
\usepackage{graphicx}
\usepackage{setspace}
\usepackage{enumerate}
\usepackage{tikz}
\usepackage{url,color}
\usepackage{xcolor}
\usepackage{chngcntr}
\usepackage{multicol}
\usepackage{comment}
\usepackage{float}
\usepackage{algpseudocode}
\usepackage{subcaption}
\usepackage[ruled]{algorithm2e}
\usepackage{booktabs}
\usepackage{caption}
\usepackage[letterpaper,margin=2cm,top=2.5cm,bottom=2.5cm]{geometry}

\AtBeginEnvironment{algorithm}{\setstretch{1.2}}

\hypersetup{colorlinks=true,citecolor=blue,linkcolor=blue,urlcolor=blue}

%%%%%%%%%%%%%%%%%%%%%%%%%%%%%%%%%%%%%%%%%%%%%%%%%%%%%%%%%%
%% Theorem environments
%%%%%%%%%%%%%%%%%%%%%%%%%%%%%%%%%%%%%%%%%%%%%%%%%%%%%%%%%%
\newtheorem{theorem}{Theorem}[section]
\newtheorem{lemma}{Lemma}[section]
\newtheorem{proposition}{Proposition}[section]
\newtheorem{corollary}{Corollary}[section]
\newtheorem{claim}{Claim}[section]
\newtheorem{conjecture}{Conjecture}[section]
\newtheorem{hypothesis}{Hypothesis}[section]
\newtheorem{assumption}{Assumption}[section]
\newtheorem{remark}{Remark}[section]
\newtheorem{example}{Example}[section]
\newtheorem{problem}{Problem}[section]
\newtheorem{definition}{Definition}[section]
\newtheorem{question}{Question}[section]
\newtheorem{condition}{Condition}[section]

%%%%%%%%%%%%%%%%%%%%%%%%%%%%%%%%%%%%%%%%%%%%%%%%%%%%%%%%%%
%% Math shorthands
%%%%%%%%%%%%%%%%%%%%%%%%%%%%%%%%%%%%%%%%%%%%%%%%%%%%%%%%%%
\newcommand{\N}{\mathbb{N}}
\newcommand{\Z}{\mathbb{Z}}
\newcommand{\E}{\mathbb{E}}
\newcommand{\Q}{\mathbb{Q}}
\newcommand{\R}{\mathbb{R}}
\newcommand{\C}{\mathbb{C}}
\newcommand{\norm}[1]{\left\lVert#1\right\rVert}

%%%%%%%%%%%%%%%%%%%%%%%%%%%%%%%%%%%%%%%%%%%%%%%%%%%%%%%%%%
%% Draft-note macros (REMOVE before submission)
%%%%%%%%%%%%%%%%%%%%%%%%%%%%%%%%%%%%%%%%%%%%%%%%%%%%%%%%%%
\newcommand{\todo}[1]{\textcolor{red}{\textbf{[TODO:}\,#1\textbf{]}}}
\newcommand{\needdata}[1]{\textcolor{orange}{\textbf{[NEED DATA:}\,#1\textbf{]}}}
\newcommand{\draft}[1]{\textcolor{blue}{#1}}
\newcommand{\note}[1]{\textcolor{purple}{\textbf{[NOTE:}\,#1\textbf{]}}}

\allowdisplaybreaks

%%%%%%%%%%%%%%%%%%%%%%%%%%%%%%%%%%%%%%%%%%%%%%%%%%%%%%%%%%
%% Title, authors, abstract
%%%%%%%%%%%%%%%%%%%%%%%%%%%%%%%%%%%%%%%%%%%%%%%%%%%%%%%%%%

\title{Refining ReCom: Kernel-Weighted Quadratic Proposals\\
       for Redistricting Markov Chains}

\author{Alaittin K\i rt\i \c{s}o\u{g}lu \\
        \small Department of Applied Mathematics, Illinois Institute of Technology \\
        \small Chicago, IL 60616 \\
        \small \texttt{kirtisoglu@hawk.iit.edu}}

\date{\today}

\begin{document}

\maketitle

\begin{abstract}
The Recombination (ReCom) algorithm and its variants dominate Markov
chain Monte Carlo (MCMC) sampling for political redistricting, but the
uniform spanning-tree split is blind to geographic structure: shared
boundary length, population density, and administrative boundaries all
carry signal that the proposal ignores, and existing weighted-tree
remedies are indicator-based heuristics with uncharacterized targets.
We propose a \emph{kernel-aware local refinement} for redistricting
MCMC, formulated as a single Laplacian energy
\(
  \min_{x}\; \alpha\, x^\top L\, x + \beta\, \|x - x_0\|^2
\)
subject to global population balance, where $L$ is an
arbitrary-kernel-weighted graph Laplacian on the merged-pair subgraph
and the fidelity term anchors each move to the current assignment.
One energy yields three regimes with honest labels: a continuous QP
relaxation with Cheeger-sweep rounding (exploratory heuristic); a
mixed-integer QP refinement inside a tempered ReCom chain (a
matheuristic ensemble generator, deliberately not a sampler); and a
QP proposal with single-shot randomized rounding and exact
Metropolis--Hastings correction, to our knowledge the first
optimization-based redistricting proposal with a detailed-balance
proof. On Iowa's 99-county dual graph ($k = 4$,
$\varepsilon = 0.05$), the integer refinement reduces mean cut edges
by $21.1\%$ versus Tempered ReCom (Cohen's $d = -2.65$,
ESS-corrected), while the corrected sampler exposes how far the
ReCom family sits below its nominal tempered target: exactly
quantifiable on the enumerated fiber, and dramatic in trend on Iowa.
A structural diagnostic of three fully enumerated partition fibers
finds genuine community structure but, correcting a common intuition,
no narrow corridors: plan space is regionally organized yet well
connected, so local chains stall because their moves are small, not
because the space is nearly disconnected, which is exactly why global
tree moves and local refinements are complementary.
We also interpret the QP-diffusion as an implicit Euler step of the
graph heat equation, exhibit its closed-form spectral filter, and
document two balance-constraint bug classes (one affecting bare
GerryChain \texttt{recom} proposals) that materially change measured
results.
\end{abstract}

\tableofcontents
\newpage

%════════════════════════════════════════════════════════════════════════════════
\section{Introduction}\label{sec:intro}
%════════════════════════════════════════════════════════════════════════════════

\subsection{The redistricting problem and ReCom's role}\label{sec:intro-problem}

Political redistricting seeks a partition of a geographic region into
electoral districts satisfying legal, demographic, and geometric
requirements: contiguity, population balance, compactness, and
preservation of communities of interest. Computational approaches model
the problem on the dual adjacency graph $G = (V,E)$ of geographic units
(precincts, census blocks, counties). A valid $k$-district plan is a
partition
\[
  P = \{V_1, \ldots, V_k\}, \qquad \bigsqcup_{i=1}^k V_i = V,
\]
where each $G[V_i]$ is connected and the populations $p(V_i)$ are
approximately equal. Let $\mathcal{P}_k(G)$ denote the space of valid
$k$-district plans. The cardinality of $\mathcal{P}_k(G)$ is enormous:
even four congressional districts over the 99 counties of Iowa admit
roughly $10^{24}$ plans~\cite{fifield2020essential}.

Markov chain Monte Carlo (MCMC) has become the standard tool for
sampling $\mathcal{P}_k(G)$. Early flip chains modify a single
precinct's assignment per step~\cite{fifield2020automated} and mix
slowly. The Recombination (ReCom) algorithm of DeFord, Duchin, and
Solomon~\cite{deford2021recombination} performs multi-precinct updates
by merging two adjacent districts, sampling a uniform random spanning
tree of the merged subgraph, and cutting an edge to recover a balanced
bipartition. ReCom dramatically improved the mixing efficiency of
redistricting MCMC and is now the workhorse of the field, with
reversible~\cite{cannon2022spanning} and multi-scale~\cite{autry2020multiscale}
extensions.

\textbf{ReCom's blind spot.} The uniform spanning-tree split treats
every edge of the merged subgraph identically: the proposed
bipartition is determined by the topology of the merged subgraph
alone. Real geography is not topology-only: shared boundary
length, population density, demographic similarity, and
administrative boundaries all carry signal that could shape the
proposed split. Practitioners have felt this: county-weighted
trees~\cite{clelland2022colorado} and GerryChain's region-aware
surcharges~\cite{gerrychain2024} inject indicator-based edge weights
into the tree draw, and SpecReCom~\cite{davies2025spectral} replaces
the tree with a Fiedler-vector cut of a randomly weighted Laplacian.
But these mechanisms are restricted (region indicators, hard-coded
county kernels, or no kernel at all), none includes a
locality-preserving fidelity term, and, as their authors
acknowledge, each shifts the chain's stationary distribution in
an uncharacterized, uncorrected way.

\subsection{Contribution}\label{sec:intro-contribution}

We introduce a kernel-aware local-refinement framework for ReCom-style
MCMC. The core object is a single Laplacian energy on the merged-pair
subgraph $H = G[V_i \cup V_j]$:
\begin{equation}\label{eq:core-energy}
  E(x) \;=\; \alpha\, x^\top L\, x \;+\; \beta\, \|x - x_0\|^2,
  \qquad
  (1 - \varepsilon)\bar p \;\le\; p^\top x \;\le\; (1+\varepsilon)\bar p,
\end{equation}
where $L$ is the kernel-weighted graph Laplacian of $H$, $x_0$ encodes
the current bipartition, and the balance constraint is anchored to the
global per-district ideal (Equation~\eqref{eq:balance-window}). The
Laplacian term $x^\top L x$ counts (or approximates) cut edges between
the two sides; the fidelity term $\|x - x_0\|^2$ penalizes deviation
from the current assignment, enforcing locality of each move. The
paper makes four contributions:

\begin{enumerate}
  \item \textbf{A kernel-aware local refinement} (\S\ref{sec:model}).
    The Laplacian $L$ accepts any nonnegative edge-weighting kernel
    as a first-class input. We instantiate the framework with
    uniform, perimeter-weighted, and density-weighted kernels. Prior
    weighted-proposal chains restrict to region-indicator surcharges
    or hard-coded county kernels and leave the induced distribution
    uncharacterized~\cite{clelland2022colorado,gerrychain2024,
    tapp2025compactness}; to our knowledge no prior MCMC
    redistricting sampler accepts an arbitrary user-supplied edge
    kernel at the proposal level \emph{while retaining an explicit
    Metropolis correction against a stated target} (contribution 2).

  \item \textbf{One energy, three solution regimes}
    (\S\ref{sec:model}--\S\ref{sec:iqp}), with honest labels:
    \begin{itemize}
      \item \textbf{Chain B (QP-diffusion)}: convex QP over
        $x \in [0,1]^{|V_H|}$, Cheeger-sweep rounding, orphan-flip
        repair, mixed 50/50 with ReCom, an exploratory heuristic
        with no stationarity claim.
      \item \textbf{Chain C (IQP local)}: the same energy over
        $x \in \{0,1\}^{|V_H|}$, solved as a convex MIQP by Gurobi
        inside a tempered ReCom chain, a \emph{matheuristic}
        ensemble generator in the lineage of optimization-based
        recombination~\cite{brous2025recom,cannon2023voting}, not a
        sampler (\S\ref{sec:chain-c-honest}).
      \item \textbf{Chain E (QP-MH)}: the same QP proposal with
        single-shot randomized rounding and an exact
        Metropolis--Hastings correction, to our knowledge the
        first optimization-based redistricting proposal with a
        detailed-balance proof (Proposition~\ref{prop:qpmh-db}).
    \end{itemize}
    Setting $\alpha = 1, \beta = 0$ with the uniform kernel recovers
    balanced cut-minimization as a special case of Chain C; the
    $\beta \to 0^+$ limit of Chain B approaches a balanced
    Fiedler-style (BalSpecReCom-like) split
    (\S\ref{sec:iqp-subsumption}).

  \item \textbf{A structural diagnostic of the partition fiber, with
    a corrected intuition} (\S\ref{sec:fiber-corridors}). On three
    fully enumerated test beds, valid-plan space carries genuine
    community structure (modularity $0.58$--$0.69$ against
    degree-preserving nulls of $\sim 0.35$), but the folk picture of
    islands joined by narrow corridors is refuted: interfaces between
    communities are tens to thousands of edges wide, nearly all
    community pairs are directly connected, and the exact-balance
    fiber's 16 literal bridges each dangle a single pendant vertex.
    A step-trace diagnostic shows pure tree-move chains crossing
    community boundaries at $14$--$16\%$ per step (near the
    community-blind baseline) while optimization-guided local
    mixtures drop to $6$--$9\%$ and the exact local chain to
    $0.3\%$: local chains stall because their moves are small, not
    because plan space is bottlenecked. (The diagnostic isolates
    locality, not the fidelity term specifically; the evidence that
    the $\beta$ anchor matters is the Iowa compactness contrast,
    where the unanchored spectral variant matches plain ReCom and
    only the anchored integer refinement improves on it.)

  \item \textbf{An honest empirical protocol and two cautionary
    findings} (\S\ref{sec:experiments}). All comparisons use
    ESS-corrected inference and per-step global balance audits. On
    Iowa counties ($k = 4$, $\varepsilon = 0.05$), Chain C reduces
    mean cut edges by $21.1\%$ versus Tempered ReCom (Cohen's
    $d = -2.65$); Chain B's previously measured gain disappears once
    a balance-constraint bug is fixed, we document the bug class,
    which also affects bare GerryChain \texttt{recom} proposals used
    without a validator. On the enumerated $5 \times 5$ fiber, the
    exact sampler Chain E exposes how far the entire ReCom family
    sits from its nominal tempered target.
\end{enumerate}

\paragraph{Positioning.} \emph{vs.\ SpecReCom}~\cite{davies2025spectral}:
ours adds the fidelity term that turns a fresh spectral re-split into
a local refinement, accepts arbitrary kernels, enforces balance
inside the optimization, and (in Chain E) carries an exact MH
correction. \emph{vs.\ weighted-tree
proposals}~\cite{clelland2022colorado,gerrychain2024}: ours accepts
arbitrary kernels rather than region indicators and pairs weighting
with a characterized target. \emph{vs.\ the matheuristic
line}~\cite{brous2025recom,cannon2023voting}: Chain C is a stochastic
chain with an explicit diversity mechanism rather than
accept-if-better local search, and we benchmark it as an optimizer.
\emph{vs.\ RevReCom, Cycle Walk, Marked Edge Walk,
BUD}~\cite{cannon2022spanning,deford2025cycle,mcwhorter2025marked,
akitaya2026bud}: those chains achieve computable transition
probabilities by lifting to spanning-forest state spaces; Chain E
achieves the same property at the partition level via a deterministic
convex program plus randomized rounding, and the two routes are
complementary, their global moves could serve as the mobile
component that our exact-but-local kernel lacks.

\paragraph{Roadmap.}\label{sec:intro-roadmap}
\S\ref{sec:related} reviews related work. \S\ref{sec:background} sets
up notation, recalls ReCom, and characterizes the partition fiber's
community macro-structure, and the absence of narrow corridors, on three test beds.
\S\ref{sec:model} formulates the kernel-flexible QP-diffusion model in
continuous (Chain B), integer (Chain C), and Metropolis-corrected
(Chain E) forms, with Cheeger-sweep rounding, orphan-flip repair,
hybrid composition with ReCom, and a brief spectral interpretation.
\S\ref{sec:experiments} reports balance-audited, ESS-corrected
empirical results on Iowa, the $5 \times 5$ enumerated fiber, and
\draft{NC precincts (pending Sprint 3)}. \S\ref{sec:discussion}
discusses roles, trade-offs, and limitations;
\S\ref{sec:conclusion} concludes. Appendices contain algorithm
details, hyperparameter tables, full sweep tables, and a
reproducibility checklist.

%════════════════════════════════════════════════════════════════════════════════
\section{Related work}\label{sec:related}
%════════════════════════════════════════════════════════════════════════════════

\paragraph{MCMC for redistricting.}\label{sec:related-mcmc}

Early MCMC for redistricting framed the problem as a graph cut and used
single-unit flip chains~\cite{fifield2020automated} or
Swendsen--Wang-style large-swap moves. Flip chains mix poorly and
produce highly irregular districts~\cite{deford2021recombination,
najt2024empirical}. ReCom~\cite{deford2021recombination} introduced
the spanning-tree merge-split move and is now standard. Extensions
include Multi-Scale Merge--Split~\cite{autry2020multiscale} preserving
nested communities (counties, precincts), Forest ReCom~\cite{autry2023metropolized},
Sequential Monte Carlo~\cite{mccartan2023sequential}, multiscale
parallel tempering~\cite{chuang2025multiscale}, and reversible
recombination~\cite{cannon2022spanning} targeting the canonical
spanning-tree distribution. A consistent theme of the 2022--2026
literature is that a sampler is only as defensible as its target:
RevReCom~\cite{cannon2022spanning}, Metropolized forest
recombination~\cite{autry2023metropolized}, the cycle
walk~\cite{deford2025cycle}, and the marked edge
walk~\cite{mcwhorter2025marked} all pay a structural price (lifted
state spaces, spanning-forest counting) precisely to keep transition
probabilities computable. Mixing-time bounds remain open for
all variants in the practical regime; Akitaya et
al.~\cite{akitaya2022reconfiguration} prove that even deciding
reachability between balanced partitions under recombination is
PSPACE-complete.

\paragraph{Weighted and geography-aware proposals.}\label{sec:related-weighted}

Edge weighting inside the ReCom proposal is not new, and we position
against it explicitly. Clelland et
al.~\cite{clelland2022colorado} replace the proposal's edge-weighting
function with \emph{county-weighted trees}, intra-county edge
weights drawn from $\mathcal{U}[0, w]$ for a tunable scalar $w$, to reduce county splits in Colorado ensembles. GerryChain's
region-aware ReCom~\cite{gerrychain2024} productionizes the same idea:
user-specified surcharges are added to the randomized Kruskal weights
for edges crossing designated region boundaries. Tapp~\cite{tapp2025compactness}
modifies the ReCom tree-and-cut step to improve shape-based
compactness at the proposal level. In all three, the weighting is an
indicator- or attribute-indexed heuristic, and (as their authors
acknowledge) the induced stationary distribution is uncharacterized
and uncorrected. Conversely, the samplers with known targets cited
above handle only unweighted spanning-tree/forest measures or other
tree-count-based families. The gap this paper addresses sits at the
intersection: a proposal family that accepts an \emph{arbitrary}
nonnegative edge kernel as a first-class input, together with (for the
Metropolis-corrected variant of \S\ref{sec:qp-mh}) an explicit
accept/reject step against a stated target.

\paragraph{Optimization-based districting and matheuristics.}\label{sec:related-opt}

A parallel OR literature treats districting as exact optimization:
the Hess model lineage, contiguity-constrained
MIPs~\cite{validi2022imposing}, and cut-edge minimization solved to
optimality at county/tract scale~\cite{validi2022political}. Between
the sampling and optimization families~\cite{becker2022redistricting},
hybrids are emerging: short bursts~\cite{cannon2023voting} use ReCom
as a stochastic optimizer; Brous and Shmoys~\cite{brous2025recom}
replace the spanning-tree re-split with a Gurobi integer program
inside an accept-if-better local search for majority-minority
districts. Our integer chain (\S\ref{sec:iqp}) is closest to this
matheuristic line (it embeds convex MIQP local refinements inside a
recombination chain) and we present it as such: an
optimization-flavored ensemble generator, not a sampler with a
stationarity guarantee (\S\ref{sec:chain-c-honest}). The proximal
term $\beta\|x - x_0\|^2$ is the quadratic analogue of local
branching~\cite{fischetti2003local}, imposing a soft trust region
around the incumbent.

\paragraph{Spectral partitioning and diffuse-interface models.}\label{sec:related-spectral}

Spectral graph partitioning has a long history in image segmentation
and clustering~\cite{shi2000normalized}. In redistricting, the recent
SpecReCom of Davies, Job, Kampbell, Kim, and
Seo~\cite{davies2025spectral} (arXiv preprint, June 2025) replaces the
random spanning tree with a Fiedler-vector sign cut on a randomly
edge-weighted Laplacian of the merged pair, reporting roughly half the
cut edges of RevReCom on a $56 \times 56$ grid and a Colorado planar
graph; plain SpecReCom sacrifices population balance, which their
BalSpecReCom variant partially restores. Diffuse-interface models
build on the Ginzburg--Landau (GL) functional, combining a smoothness
term with a fidelity term; the Chan--Vese model~\cite{chan2001active}
and Merriman--Bence--Osher (MBO) threshold dynamics implement this
iteratively, and Bertozzi and collaborators~\cite{bertozzi2012diffuse,
merkurjev2013mbo} port the machinery to weighted graphs, where the
smoothness term becomes the Laplacian quadratic form $x^\top L x$ and
graph Tikhonov regularization~\cite{pilavci2021graph,shuman2013emerging}
recovers heat-kernel smoothing. Jacobs and Walch~\cite{jacobs2018pde}
brought threshold dynamics to redistricting before SpecReCom, using
volume-constrained MBO via auction dynamics~\cite{jacobs2018auction}
to generate compact equal-population maps, a continuum PDE
approach, not a graph-kernel one, and not embedded in MCMC. To our
knowledge the present work is the first to use the \emph{graph} GL
lineage (kernel-weighted Laplacian smoothing plus fidelity) as a
proposal inside redistricting MCMC.

\paragraph{Cut minimization and discrete compactness.}\label{sec:related-cut}

The minimum graph cut~\cite{stoer1997simple,karger2000minimum} and its
balanced variants such as normalized cut~\cite{shi2000normalized} are
classical objectives. NP-hardness of balanced minimum
cut~\cite{wagner1993between} motivates continuous relaxations, of
which the Fiedler-vector cut and the Total Variation Isoperimetric
Profile~\cite{deford2019total} are two well-known examples. The cut
edge count $|\partial P|$ has been proposed as a discrete compactness
measure for redistricting~\cite{duchin2018discrete} and shown to
correlate well with classical geometric scores (Polsby--Popper, Reock)
while being more stable under unit refinement.

%════════════════════════════════════════════════════════════════════════════════
\section{Background: ReCom and the partition fiber}\label{sec:background}
%════════════════════════════════════════════════════════════════════════════════

\subsection{Dual graph and \texorpdfstring{$k$}{k}-district plans}\label{sec:background-graph}

Let $G = (V, E)$ be the dual graph of a state, with vertices
representing geographic units (precincts, census blocks, counties) and
edges encoding shared boundaries. Each vertex $v$ carries a population
$p_v \ge 0$ and may carry additional attributes (demographic vectors,
shapefile geometry). A \emph{$k$-district plan} is a partition
$P = \{V_1, \ldots, V_k\}$ of $V$ such that each induced subgraph
$G[V_i]$ is connected (\emph{contiguity}) and the populations satisfy
\begin{equation}\label{eq:pop-balance}
  (1 - \varepsilon)\bar p \;\le\; p(V_i) \;\le\; (1 + \varepsilon)\bar p,
  \qquad \bar p := \frac{p(V)}{k},
\end{equation}
for a tolerance $\varepsilon \in [0, 1)$. We write $\mathcal{P}_k(G; \varepsilon)$
for the set of valid plans. The \emph{cut edges} of a plan are
\[
  \partial P \;=\; \{(u,v) \in E : P(u) \ne P(v)\},
\]
and we call $|\partial P|$ the \emph{cut-edge count} of $P$. This
quantity is the discrete analogue of the perimeter of the district
boundaries and serves as our primary compactness
metric~\cite{duchin2018discrete}.

\begin{figure}[H]
  \centering
  \IfFileExists{fig/redistricting.png}{%
    \includegraphics[width=0.7\textwidth]{fig/redistricting.png}%
  }{%
    \fbox{\begin{minipage}{0.7\textwidth}\vspace{2.5cm}\centering
      \ttfamily\small\textcolor{orange}{[Missing figure:
      \texttt{fig/redistricting.png}]}\vspace{2.5cm}\end{minipage}}%
  }
  \caption{The dual graph of Iowa's 99 counties (left) and the four
    congressional districts enacted after the 2010 census (right).
    Nodes sit at county centroids; two counties are joined by an edge
    when they share a boundary segment.}
  \label{fig:iowa-dual}
\end{figure}

\subsection{Recombination (ReCom)}\label{sec:background-recom}

Recombination~\cite{deford2021recombination} is a Markov chain on
$\mathcal{P}_k(G; \varepsilon)$ whose single step is:

\begin{algorithm}[H]
\caption{Vanilla ReCom step (Chain V)}\label{alg:recom}
\KwIn{Current plan $P$, dual graph $G$, balance tolerance $\varepsilon$.}
\KwOut{Next plan $P'$.}
Select two adjacent districts $V_i, V_j$ uniformly at random from boundary pairs;\\
Form the induced subgraph $H = G[V_i \cup V_j]$;\\
\Repeat{a cut $e^*$ exists yielding two $\varepsilon$-balanced components}{
  Sample a uniform random spanning tree $T$ of $H$ via Wilson's algorithm~\cite{wilson1996generating};\\
  Search edges of $T$ for a cut $e^*$ whose removal produces two $\varepsilon$-balanced components.
}
Define $V_i', V_j'$ as the two components of $T \setminus \{e^*\}$;\\
$P' \gets P$ with $V_i, V_j$ replaced by $V_i', V_j'$;\\
\Return $P'$.
\end{algorithm}

\noindent The tree-based proposal flips many precincts simultaneously
between $V_i$ and $V_j$ while preserving contiguity automatically (any
edge removal from a tree produces two connected components). This is
the key advantage over single-precinct flip chains, which traverse the
state space too slowly to be practically useful.

\textbf{Chains compared.} We compare six chains, three of them from
this work:

\begin{table}[H]
\centering
\small
\begin{tabular}{llll}
\toprule
Chain & Method & Mechanism & Stationarity claim \\
\midrule
V & Vanilla ReCom~\cite{deford2021recombination} & spanning-tree split, always accept & unknown target \\
A & Tempered ReCom & ReCom + MH on cuts (symmetric-$q$) & approximate \\
B & ReCom + QP-diffusion mix (this work) & continuous QP, Cheeger sweep, repair & none (heuristic) \\
C & Tempered ReCom + IQP local (this work) & MIQP refinement, orphan-flip repair & none (matheuristic) \\
D & ReCom + spectral mix (after~\cite{davies2025spectral}) & Fiedler + balanced sweep, no $\beta$, no repair & none (heuristic) \\
E & QP-MH (this work, \S\ref{sec:qp-mh}) & QP + Bernoulli rounding + exact MH & \textbf{exact} \\
\bottomrule
\end{tabular}
\caption{The six chains compared in this paper. The last column is the
  chain's relationship to the tempered target
  $\pi(P) \propto \exp(-\lambda|\partial P|)$: only Chain E carries an
  exact Metropolis--Hastings correction; Chains B--D are heuristics
  whose output distributions we characterize empirically, and Chain A
  relies on the symmetric-$q$ approximation
  (\S\ref{sec:qp-mh}).}
\label{tab:five-chains}
\end{table}

Tempered ReCom uses a Metropolis--Hastings filter targeting
$\pi(P) \propto \exp(-\lambda |\partial P|)$ with the symmetric-$q$
approximation; see Appendix~\ref{app:hyperparams} for the locked-in
$\lambda = 0.10$. Our Chain D departs from published SpecReCom in one
deliberate way: instead of the sign split of the Fiedler vector (which
ignores population balance), it runs our balance-constrained Cheeger
sweep over the Fiedler order, closer in spirit to BalSpecReCom, so that all chains produce valid plans and the comparison isolates the
role of the fidelity term. All chains additionally enforce the global
two-sided balance window of Equation~\eqref{eq:balance-window} on
every move; we note (\S\ref{sec:exp-balance}) that the raw GerryChain
\texttt{recom} proposal alone does not.

\paragraph{ReCom's two weaknesses.}\label{sec:background-weakness}
ReCom is empirically strong but has two known weaknesses that motivate
the present work.

\textbf{(W1) Non-compact splits.} Uniform spanning trees of the
merged subgraph favor topologically irregular partitions because the
number of irregular subtrees is combinatorially much larger than the
number of compact ones~\cite{deford2021recombination,deford2019total}.
The cut edge is chosen \emph{after} the tree is drawn, so the
proposal cannot anticipate compactness.

\textbf{(W2) Kernel-blindness.} The uniform spanning tree weights
every edge equally. Shared-boundary length, population density,
demographic similarity, and administrative boundary attributes are
ignored at the proposal level. They can only enter via post-hoc
filtering, which raises rejection rates and is hard to tune.

The diffusion-based proposal of \S\ref{sec:model} addresses both
weaknesses by replacing the spanning-tree-cut step with a single
convex optimization that minimizes a kernel-weighted Laplacian energy
with a locality-preserving fidelity term. The energy
$x^\top L\, x = \sum_{(u,v) \in E} w_{uv} (x_u - x_v)^2$ equals the
weighted cut count for binary $x$, so minimizing it directly minimizes
the discrete perimeter.

\subsection{Partition fiber structure: communities without bottlenecks}\label{sec:fiber-corridors}

This subsection documents a structural diagnostic of the partition fiber
$\mathcal{P}_k(G; \varepsilon)$, viewed as a graph under its minimal
valid move relation. Two findings emerge, one positive and one
negative, and we report both. Positively, all three enumerated fibers
carry genuine mesoscale community structure, with modularity far above
degree-preserving null models, and chains differ sharply in how often
they cross community boundaries. Negatively, a plausible folk
intuition (that the communities are ``islands'' joined by narrow
corridors or bottleneck edges) is refuted by method-independent
statistics: community interfaces are wide, nearly all community pairs
are directly connected, and the literal bridges of the exact-balance
fiber turn out to be pendant vertices. We had entertained the
islands-and-corridors reading in an earlier draft of this work; the
measurements below correct it. The diagnostic still motivates the
local-refinement design of \S\ref{sec:model}, and we believe both
findings are of independent interest for analyzing redistricting
MCMC.

\paragraph{Three test beds.}\label{sec:fiber-testbeds}

We construct three fully enumerated partition fibers spanning grid
size, district count, and balance tolerance:

\begin{table}[H]
\centering
\small
\begin{tabular}{lrrrrrl}
\toprule
Fiber & Nodes & Edges & $k$ & $\varepsilon$ & Sizes & Adjacency \\
\midrule
$5 \times 5$, $\varepsilon = 0$    & 4\,006   & 13\,416  & 5 & 0    & \{5\}       & boundary swap \\
$5 \times 5$, $\varepsilon = 0.20$ & 193\,128 & 969\,760 & 5 & 0.20 & \{4, 5, 6\} & single flip \\
$4 \times 4$, $\varepsilon = 0.25$ & 1\,953   & 7\,024   & 4 & 0.25 & \{3, 4, 5\} & single flip \\
\bottomrule
\end{tabular}
\caption{The three test-bed fibers. Adjacency is the \emph{minimal
  valid move} of each fiber: at $\varepsilon = 0$ no single-vertex
  reassignment preserves exact balance (it produces districts of sizes
  4 and 6), so adjacency is the boundary \emph{swap}: exchange one
  boundary vertex of district $a$ with one of district $b$, both
  districts remaining connected. At $\varepsilon > 0$ adjacency is the
  single \emph{flip}: reassign one boundary vertex, both affected
  districts remaining connected and within size bounds. Cross-fiber
  comparisons must be read with this caveat: the $\varepsilon = 0$
  fiber lives under a different (two-vertex) move relation than the
  $\varepsilon > 0$ fibers.}\label{tab:fibers}
\end{table}

For each fiber, all valid partitions are enumerated and the single-flip
adjacency graph is built explicitly (enumeration runtimes on a 2024
MacBook range from $\sim 5$ s for the $4 \times 4$,
$\varepsilon = 0.25$ fiber to $\sim 50$ min for the $5 \times 5$,
$\varepsilon = 0.20$ fiber; scripts accompany the replication
materials).

\paragraph{Community structure, validated against null models.}\label{sec:fiber-islands}

Each fiber graph carries substantial mesoscale community structure
under modularity maximization. We report the method per fiber (greedy
modularity, Clauset--Newman--Moore, for the $\varepsilon = 0$ and
$4 \times 4$ fibers; Louvain, resolution $1.0$, seed 42, for the large
$\varepsilon = 0.20$ fiber), validate against degree-preserving
rewiring nulls, and treat the specific community counts as
method-dependent descriptive statistics:

\begin{table}[H]
\centering
\small
\begin{tabular}{lrrrrl}
\toprule
Fiber & Communities & $Q$ & Null $Q$ & Community sizes & Method \\
\midrule
$5 \times 5$, $\varepsilon = 0$    & 10 & 0.65 & 0.36 & 4--788 (six $\ge$ 579) & CNM \\
$5 \times 5$, $\varepsilon = 0.20$ & 26 & 0.69 & n/a  & $\sim 2\,800$--$13\,350$ & Louvain \\
$4 \times 4$, $\varepsilon = 0.25$ & 7  & 0.58 & 0.34 & 33--475 & CNM \\
\bottomrule
\end{tabular}
\caption{Community decompositions with degree-preserving null
  baselines (mean $Q$ of CNM on double-edge-swap rewirings of the same
  graph; the $193$k-node fiber's null is omitted for cost). Real $Q$
  exceeds null $Q$ by $0.24$--$0.29$: the community structure is
  genuine, not a degree-sequence artifact. Counts vary with the
  detection method (an earlier analysis of the $4 \times 4$ fiber with
  different settings reported 16 communities) and should be read
  qualitatively.}
\label{tab:communities}
\end{table}

\paragraph{No narrow corridors: a corrected intuition.}\label{sec:fiber-corridors-invariant}

It is tempting to picture these communities as islands joined by a few
narrow corridors, with bottleneck edges controlling passage; an
earlier draft of this work did. Method-independent statistics refute
that picture on all three fibers:

\begin{table}[H]
\centering
\small
\begin{tabular}{lrrrr}
\toprule
Fiber & Pairs connected & Inter-comm.\ edges & Interface width (med/max) & Bridges \\
\midrule
$5 \times 5$, $\varepsilon = 0$    & 15/15 (large pairs) & 18.3\% & 142 / 383 & 16, all pendant \\
$5 \times 5$, $\varepsilon = 0.20$ & 320/325 & 26.8\% & 667 / 3\,006 & 0 \\
$4 \times 4$, $\varepsilon = 0.25$ & 21/21   & 23.3\% & 71 / 186 & 0 \\
\bottomrule
\end{tabular}
\caption{Anti-bottleneck statistics. Nearly every pair of communities
  shares a direct interface; interfaces are tens to thousands of edges
  wide; a fifth to a quarter of all fiber edges run between
  communities; and the 16 literal single-edge bridges of the
  exact-balance fiber each split off a \emph{single vertex} (pendant
  leaves), not an island. Conductance of the six large
  $\varepsilon = 0$ communities is $0.15$--$0.22$, an order of
  magnitude above bottleneck territory.}
\label{tab:narrow-throats}
\end{table}

Concretely: on the $\varepsilon = 0.20$ fiber, interfaces thinner than
50 edges do exist (there are 16 of them), but they are simply the tail
of a width distribution whose median is 667 and whose top interfaces
carry $3\,000$ edges; a chain moving through this graph crosses
between communities through the wide interfaces, and the thin tail is
irrelevant to mixing. The correct summary of fiber macro-structure is
therefore \emph{communities without bottlenecks}: valid-plan space is
regionally organized but well connected, and slow traversal by local
chains (below) is a property of the \emph{dynamics}, not of corridor
topology. This negative result matters for chain design: it says
local chains stall not because the state space is nearly disconnected
but because their moves are small and, in Chain E's case, because
distant proposals are rarely accepted; global recombination moves, by
contrast, relocate across communities in a single step.

\paragraph{Chain step-trace diagnostic.}\label{sec:fiber-trace}

The island/corridor decomposition supports a direct diagnostic for any
redistricting chain. We label each visited partition by its
greedy-modularity community on the $\varepsilon = 0$ fiber and record
the \emph{jump rate}, the fraction of steps that move to a
different community.

\begin{table}[H]
\centering
\small
\begin{tabular}{lrrr}
\toprule
Chain & Communities visited & Jump rate & Jumps in 1\,000 steps \\
\midrule
V: Vanilla ReCom        & 8/10 & \textbf{15.8\%} & 158 \\
A: Tempered ReCom       & 7/10 & \textbf{13.9\%} & 139 \\
B: QP-diffusion (50/50) & 7/10 & \textbf{6.0\%}  & 60  \\
C: ReCom + IQP local    & 6/10 & 8.7\%           & 87  \\
D: ReCom + spectral mix & 7/10 & \textbf{8.2\%}  & 82  \\
E: QP-MH exact          & 2/10 & 0.3\%           & 3   \\
\bottomrule
\end{tabular}
\caption{Inter-community jump rates on the $5 \times 5$,
  $\varepsilon = 0$ fiber (1\,000 steps per chain, seed 42, per-chain
  RNG streams). A community-blind baseline crosses at roughly the
  inter-community edge share (18.3\%,
  Table~\ref{tab:narrow-throats}); ReCom variants (V, A) sit near it
  at $\sim 14$--16\% per step, the local kernels (B, C, D) are
  stickier at $\sim 6$--9\%, and the strictly local QP-MH chain (E)
  almost never crosses, a property of its small accepted moves rather
  than of any bottleneck in the fiber.}\label{tab:island-trace}
\end{table}

\begin{figure}[H]
  \centering
  \includegraphics[width=\textwidth]{../diffusion_model/recom/iowa_hybrid_test_out/island_trace_5x5.png}
  \caption{Chain step-traces over $5 \times 5$ communities. Each panel
    shows the sequence of community labels visited by one chain over
    1\,000 steps; horizontal bands are communities, x-axis is time.
    Vanilla and Tempered ReCom oscillate rapidly between communities
    (high jump rate); the anchored local kernels dwell within
    communities for extended runs, and the exact QP-MH chain (bottom)
    barely leaves its starting community.}
  \label{fig:island-trace}
\end{figure}

\textbf{Reading of the diagnostic.} ReCom relocates across
communities in single steps; optimization-guided local kernels work
within them; and the exact-but-strictly-local Chain E shows the
failure mode of small accepted moves. The regimes are complementary,
motivating the hybrid chain design of \S\ref{sec:model}. We note the
diagnostic's limits: one run per chain on one fiber, and community
labels from a specific detection method, we read it as motivating evidence, not proof.

\paragraph{Implication for chain design.}\label{sec:fiber-implication}

The diagnostic suggests an explicit principle:

\begin{quote}
\itshape A good redistricting sampler should combine global moves
that relocate across regions of plan space with local moves that
refine within them. ReCom does the former; anchored local kernels do
the latter.
\end{quote}

This is the design principle of Chains B and C: each step is a random
mixture of a ReCom (across-community) move and a diffusion (within-community)
move. The remainder of the paper develops and validates this design.

%════════════════════════════════════════════════════════════════════════════════
\section{The QP-diffusion model: continuous and integer formulations}\label{sec:model}
%════════════════════════════════════════════════════════════════════════════════

\subsection{Locality kernels}\label{sec:kernels}

Choosing the edge weights $w_{ij}$ is the most consequential modeling
decision. The kernel determines what \emph{smoothness} means and
therefore controls how district boundaries move under the relaxation.
We restrict all weights to geographic neighbors,
\begin{equation}\label{eq:kernel-local}
  w_{ij} = 0 \quad \text{if } (v_i, v_j) \notin E,
\end{equation}
ensuring that diffusion operates only through geographic adjacency and
preventing long-range interactions that could produce disconnected
districts.

\paragraph{Uniform kernel ($k_{\mathrm{unif}}$).} $w_{ij} = 1$ for
every edge. The Laplacian $L$ is the combinatorial graph Laplacian;
$x^\top L x$ is exactly the cut edge count. In the $\beta \to 0^+$
limit the proposal approaches a balanced Fiedler-style split
(\S\ref{sec:iqp-subsumption}).

\paragraph{Perimeter kernel ($k_{\mathrm{perim}}$).} $w_{ij} = |\partial v_i \cap \partial v_j|$
(shared-boundary length). The Laplacian quadratic form $x^\top L x$
is then the \emph{geographic perimeter} of the partition, weighted by
boundary segment length. This kernel discourages fragmented districts
and favors boundaries that cut along short segments.

\paragraph{Density kernel ($k_{\mathrm{dens}}$).} Let $\rho_v$ denote
the population density at vertex $v$ (population per unit area). A
Gaussian density kernel
\[
  w_{ij} \;=\; |\partial v_i \cap \partial v_j| \cdot
  \exp\!\Bigl( -\frac{(\rho_i - \rho_j)^2}{\sigma^2} \Bigr)
\]
encourages cuts along sharp density gradients (urban--rural
boundaries) while preserving the topological role of shared boundary
length.

\paragraph{Demographic kernel ($k_{\mathrm{demog}}$).} Let
$d_v \in \R^m$ be a vector of demographic attributes. A Gaussian
similarity kernel $w_{ij} = \exp(-\|d_i - d_j\|^2 / \sigma^2)$
encourages communities of interest. We define this kernel for
completeness but do not deploy it in our experiments, and a caution
stronger than ``use as a soft regularizer'' is warranted: if the
attribute vector encodes race, a proposal that systematically shapes
district boundaries around racial similarity risks making race the
predominant districting criterion, which triggers strict scrutiny
under the \emph{Shaw v.\ Reno} line of racial-gerrymandering cases.
Any practical use of demographic kernels should be reviewed against
that doctrine, should document the kernel's weight relative to
race-neutral criteria, and is more defensible for attributes that are
legitimate districting factors in the jurisdiction (e.g., municipal
membership, communities of interest as defined by state law).

\paragraph{Composite kernels.} Linear combinations $w_{ij} = \sum_k a_k w_{ij}^{(k)}$
with $a_k \ge 0$ remain valid (preserve nonnegativity and locality).

\subsection{The convex QP}\label{sec:qp-model}

Let $H = G[V_i \cup V_j]$ be the merged subgraph of two adjacent
districts. Let $L_H$ be the kernel-weighted Laplacian of $H$, let
$x_0 \in \{0,1\}^{|V_H|}$ encode the current bipartition (1 for
vertices in $V_i$, 0 for vertices in $V_j$), and let $p$ be the
population vector restricted to $H$. Fix weights $\alpha, \beta > 0$
and balance tolerance $\varepsilon \in (0, 1)$. Crucially, the balance
constraint is anchored to the \emph{global} per-district ideal
$\bar p = p(V)/k$ from Equation~\eqref{eq:pop-balance}, not to
$p(V_H)/2$: both sides of the split must individually satisfy the
plan-validity bound, which yields the two-sided window
\begin{equation}\label{eq:balance-window}
  \ell \;=\; \max\bigl((1-\varepsilon)\bar p,\; p(V_H) - (1+\varepsilon)\bar p\bigr),
  \qquad
  u \;=\; \min\bigl((1+\varepsilon)\bar p,\; p(V_H) - (1-\varepsilon)\bar p\bigr).
\end{equation}
Anchoring to $p(V_H)/2$ instead (a tempting simplification, and one
present in an earlier version of our own implementation) only bounds
the local split: because $p(V_H)$ itself drifts as the chain moves
population between pairs, district populations can ratchet outside the
global tolerance by a factor of up to $(1+\varepsilon)$ per move on
heterogeneous graphs. Our experiments audit the global bound at every
step (\S\ref{sec:experiments}).

\begin{definition}[Continuous QP-diffusion]\label{def:qp}
The continuous QP-diffusion bipartition of $H$ is
\begin{align}
  x^\star \;=\; \arg\min_{x \in [0,1]^{|V_H|}} \;\;
    & \alpha\, x^\top L_H\, x \;+\; \beta\, \|x - x_0\|_2^2 \label{eq:qp-obj} \\
  \text{subject to} \;\;
    & \ell \;\le\; p^\top x \;\le\; u,
    \label{eq:qp-balance}
\end{align}
with $\ell, u$ as in Equation~\eqref{eq:balance-window}.
\end{definition}

The first term in~\eqref{eq:qp-obj} is the kernel-weighted Laplacian
smoothness energy: for binary $x$ it equals the weighted cut count
$\sum_{(u,v) \in E(H)} w_{uv}\,\mathbf{1}[x_u \ne x_v]$, and the
continuous relaxation $x \in [0,1]^{|V_H|}$ extends this to a smooth
objective on the unit cube. The second term is the \emph{fidelity term},
which penalizes deviation from the current bipartition and ensures
that each proposed move is local. The balance constraint enforces
$\varepsilon$-balance on the proposed split.

\begin{proposition}[Convexity and well-posedness]\label{prop:convex}
The objective in~\eqref{eq:qp-obj} is strictly convex on
$[0,1]^{|V_H|}$ (since $\beta > 0$ and $L_H$ is positive semidefinite).
The feasible region defined by~\eqref{eq:qp-balance} is a non-empty
polytope whenever the current plan is valid (it then contains $x_0$).
Hence the QP admits a unique minimizer $x^\star$, computable in
polynomial time by any convex QP solver.
\end{proposition}

\paragraph{Spectral interpretation.} Projecting the interior optimality
condition $(\alpha L_H + \beta I) x^\star = \beta x_0$ onto the
eigenbasis of $L_H$ with eigenvalues
$0 = \lambda_1 \le \lambda_2 \le \cdots \le \lambda_n$ yields the
closed-form frequency response
\begin{equation}\label{eq:freq-response}
  \hat x^\star_k \;=\; \frac{\beta}{\alpha \lambda_k + \beta}\,
  \hat x_{0,k},
  \qquad \hat y_k := \phi_k^\top y.
\end{equation}
The QP therefore acts as a low-pass spectral filter on the prior $x_0$,
with cutoff controlled by the ratio $\alpha / \beta$. Equivalently, with
$\tau := \alpha / \beta$, the optimum satisfies
$(x^\star - x_0)/\tau = -L_H x^\star$, the implicit Euler
discretization of the graph heat equation $\dot x = -L_H x$. Inspecting
the leading eigenvectors $\phi_2, \phi_3, \ldots$ provides a sanity
check on kernel choice: if they look like noise rather than coherent
geographic regions, the chosen kernel is probably not informative.

\subsection{Cheeger-sweep rounding and orphan-flip repair}\label{sec:rounding}

The continuous optimum $x^\star \in [0,1]^{|V_H|}$ must be rounded to a
binary bipartition. We use a \emph{Cheeger sweep}~\cite{kelner2009spectral}:
sort the entries of $x^\star$ in descending order, and for each
threshold $k = 1, \ldots, |V_H| - 1$ form the candidate set $S_k$ of the
top-$k$ vertices. Among candidates that satisfy $\varepsilon$-balance,
return the one minimizing the cut count
$|\partial(S_k, V_H \setminus S_k)|$. Candidates with disconnected sides
are passed through the orphan-flip repair below before evaluation.

\begin{algorithm}[H]
\caption{Cheeger-sweep rounding with orphan-flip repair}\label{alg:cheeger-sweep}
\KwIn{Continuous optimum $x^\star \in [0,1]^{|V_H|}$, population
  vector $p$, target $\bar p$, tolerance $\varepsilon$, merged subgraph $H$.}
\KwOut{Binary bipartition $S \subset V_H$ or \texttt{FAIL}.}
Sort $V_H$ by $x^\star$ descending; call the order $v_{(1)}, \ldots, v_{(|V_H|)}$;\\
$\mathcal{C} \gets \emptyset$ \tcp*{candidates}
\For{$k = 1$ \KwTo $|V_H| - 1$}{
  $S_k \gets \{v_{(1)}, \ldots, v_{(k)}\}$;\\
  \uIf{$p(S_k) \in [(1-\varepsilon)\bar p, (1+\varepsilon)\bar p]$ and both sides connected}{
    Add $S_k$ to $\mathcal{C}$ with cut $|\partial(S_k, V_H \setminus S_k)|$;
  }\ElseIf{$p(S_k)$ in balance and sides have $> 1$ components}{
    $S_k' \gets$ orphan-flip repair of $S_k$ (Algorithm~\ref{alg:orphan-flip});\\
    \If{$S_k'$ satisfies balance and connectivity}{
      Add $S_k'$ to $\mathcal{C}$;
    }
  }
}
\lIf{$\mathcal{C} = \emptyset$}{\Return \texttt{FAIL}}
\Return $\arg\min_{S \in \mathcal{C}} |\partial(S, V_H \setminus S)|$.
\end{algorithm}

The sweep can return a candidate $S_k$ whose induced subgraph $G[S_k]$
or $G[V_H \setminus S_k]$ is disconnected, a single boundary vertex
on the wrong side of the threshold is enough. Discarding such
candidates outright would inflate the rejection rate; instead we repair
them. The \emph{orphan-flip} routine identifies the disconnected
components of each side, keeps the largest component, and flips every
smaller (\emph{orphan}) component to the opposite side. Population
balance is then re-verified; if the repair preserves both balance and
connectivity, the repaired candidate enters the Cheeger sweep's
candidate pool. Algorithm~\ref{alg:orphan-flip} formalizes this:

\begin{algorithm}[H]
\caption{Orphan-flip repair}\label{alg:orphan-flip}
\KwIn{Candidate bipartition $(S, V_H \setminus S)$.}
\KwOut{Repaired bipartition or \texttt{FAIL}.}
Identify connected components of $G[S]$ and $G[V_H \setminus S]$;\\
\For{each side}{
  Keep the largest connected component;\\
  Flip every smaller component to the opposite side;
}
Verify both final sides are connected and $\varepsilon$-balanced;\\
\lIf{verification fails}{\Return \texttt{FAIL}}
\Return repaired $(S', V_H \setminus S')$.
\end{algorithm}

\textbf{Why repair matters.} On heterogeneous-population graphs such
as Iowa, the unrepaired Cheeger sweep produces near-feasible
candidates whose few orphan vertices (often singletons at the
boundary) can be reassigned without affecting the cut count.
Without repair, post-rounding rejection rates on Iowa run at 5--15\%;
with repair, $<3\%$. SpecReCom~\cite{davies2025spectral} does not use
repair (and we keep it out of our faithful Chain D
implementation); without the $\beta$ fidelity term, repair tends to
accept lower-quality cuts.

\paragraph{On sweep guarantees.} The classical Cheeger
inequality~\cite{kelner2009spectral} bounds the conductance of the
sweep cut of the \emph{Fiedler vector}:
$\Phi(H) \le \sqrt{2 \lambda_2(L_H^{\mathrm{sym}})}$. We deliberately
do not claim this bound for our rounded solution: the sweep here is
applied to the constrained proximal optimum $x^\star$ of
Definition~\ref{def:qp}, not to $\phi_2$, and the population-balance
restriction on admissible thresholds voids the standard argument. The
sweep should be read as a rounding heuristic whose output quality we
measure empirically; when $\alpha/\beta \to \infty$ and the balance
window is inactive, $x^\star$ tends toward the low-frequency eigenspace
and the classical intuition is recovered, but we leave a quantitative
statement to future work. A second caveat: the sweep selects the
candidate minimizing the \emph{unweighted} cut $|\partial(S_k, V_H
\setminus S_k)|$ (the compactness metric of the paper) even when
the QP energy uses a non-uniform kernel. Under non-uniform kernels the
two objectives diverge, and a kernel-faithful variant would score
candidates by the weighted cut; we flag this as a configuration choice
in the multi-kernel experiments (\S\ref{sec:exp-multikernel}).

\subsection{Hybrid composition with ReCom (Chain B)}\label{sec:hybrid-mh}

Chain B's single step is a random mixture: with probability
$p_R = 0.5$, take a Vanilla ReCom step (Algorithm~\ref{alg:recom});
otherwise take a QP-diffusion step (Definition~\ref{def:qp} +
Algorithm~\ref{alg:cheeger-sweep}). We state plainly what Chain B is
and is not: \emph{neither branch carries a Metropolis correction}, the ReCom branch always accepts (as in vanilla ReCom), and the
QP-diffusion branch applies the rounded, repaired split
unconditionally. Chain B therefore has no stationarity claim beyond
ergodicity; it is an exploratory heuristic whose empirical output
distribution we characterize against enumerated ground truth in
\S\ref{sec:exp-5x5}. The exactly corrected sampler built from the same
proposal is Chain E (\S\ref{sec:qp-mh}).

\begin{proposition}[Ergodicity of the hybrid]\label{prop:ergodic}
If ReCom is ergodic on $\mathcal{P}_k(G; \varepsilon)$ (which holds
under standard assumptions; see~\cite{deford2021recombination,
cannon2024irreducibility}), then the mixture of ReCom and
the QP-diffusion step with any $p_R \in (0, 1)$ is also ergodic.
\end{proposition}

\begin{proof}
Write the mixture kernel as $K = p_R K_{\mathrm{R}} +
(1 - p_R) K_{\mathrm{Q}}$. If $K_{\mathrm{R}}^{T}(P, P') > 0$ for
some $T$, then $K^{T}(P, P') \ge p_R^{T}\, K_{\mathrm{R}}^{T}(P, P')
> 0$, so irreducibility of ReCom implies irreducibility of the
mixture. Aperiodicity: the QP-diffusion step returns the current
plan with positive probability (the event that the rounded, repaired
split reproduces the current split of the selected pair, which
includes the positive-probability event that rounding recovers
$x_0$), so $K(P, P) > 0$ for every $P$ and the chain is aperiodic.
\end{proof}

\subsection{The exactly reversible QP-MH chain (Chain E)}\label{sec:qp-mh}

The deterministic rounding of \S\ref{sec:rounding} is what blocks a
Metropolis correction: an argmin-plus-sweep map has zero reverse
proposal probability, so the Hastings ratio degenerates. Replacing the
sweep with \emph{single-shot randomized rounding} restores a tractable,
everywhere-positive proposal density and yields an exactly reversible
chain. One step of Chain E:

\begin{enumerate}
  \item Pick an adjacent district pair $(V_i, V_j)$ uniformly at random
    from the $|\mathrm{bp}(P)|$ boundary pairs of the current plan $P$;
  \item Solve the (deterministic) QP of Definition~\ref{def:qp} on
    $H$ for $x^\star \in [0,1]^{|V_H|}$;
  \item Draw the proposal $x' \in \{0,1\}^{|V_H|}$ by independent
    Bernoulli rounding from the \emph{clipped} optimum,
    $x'_v \sim \mathrm{Bern}\bigl(\mathrm{clip}(x^\star_v;
    \delta, 1-\delta)\bigr)$ with $\delta = 10^{-12}$,
    \emph{once}, no retry on infeasibility. Clipping is part of the
    proposal's definition: the QP optimum can have active box bounds
    (coordinates exactly $0$ or $1$, e.g.\ at interior vertices pinned
    by the fidelity term), and sampling from the same clipped vector
    whose density enters the acceptance ratio is what makes the
    proposal law and its stated density identical, with every binary
    $x'$ receiving positive probability;
  \item Accept with probability
    \begin{equation}\label{eq:qpmh-accept}
      \min\Bigl(1,\;
      e^{-\lambda(|\partial P'| - |\partial P|)}\,
      \frac{q(x_0 \mid P', \{i,j\})}{q(x' \mid P, \{i,j\})}\,
      \frac{|\mathrm{bp}(P)|}{|\mathrm{bp}(P')|}\Bigr),
    \end{equation}
    where $q(\cdot \mid P, \{i,j\})$ is the Bernoulli product density
    induced by the QP optimum at the corresponding state, and proposals
    that violate contiguity or the balance
    window~\eqref{eq:balance-window} are rejected outright.
\end{enumerate}

\begin{proposition}[Detailed balance of Chain E]\label{prop:qpmh-db}
Chain E satisfies detailed balance with respect to
$\pi(P) \propto \exp(-\lambda |\partial P|)$ restricted to
$\mathcal{P}_k(G; \varepsilon)$.
\end{proposition}

Throughout, fix the deterministic edge kernel $w$ (weight jitter is
never used in this chain), write $\Omega = \mathcal{P}_k(G;
\varepsilon)$ for the valid space, and for a plan $P$ and an adjacent
pair $\{i,j\} \in \mathrm{bp}(P)$ write $H = V_i \cup V_j$ and let
$x_0(P) \in \{0,1\}^{V_H}$ denote the indicator of $V_i$ on $H$ (the
pair is drawn from a sorted list, so which district plays the role of
$i$ is a deterministic function of the unordered pair). Say that
$P' \in \Omega$ is \emph{pair-adjacent} to $P \in \Omega$ if
$P' \ne P$, the two plans differ on exactly two districts $\{i,j\}$,
and $H = V_i \cup V_j$ is connected. This relation is manifestly
symmetric in $P$ and $P'$.

\begin{lemma}[Proposal regularity]\label{lem:qpmh-regular}
Let $P, P' \in \Omega$ be pair-adjacent via districts $\{i,j\}$, and
suppose $\beta > 0$. Then:
\begin{enumerate}
  \item[(i)] $\{i,j\}$ is the unique unordered pair whose selection
    can generate the move $P \to P'$ or the move $P' \to P$, and it
    is selectable from both endpoints:
    $\{i,j\} \in \mathrm{bp}(P) \cap \mathrm{bp}(P')$.
  \item[(ii)] The QP of Definition~\ref{def:qp} on $H$ with the
    global window~\eqref{eq:balance-window} has the same (nonempty)
    feasible set for the forward and reverse solves, and its optimum
    is unique, hence a deterministic map of the pair
    $(\text{endpoint}, \{i,j\})$.
  \item[(iii)] The clipped-Bernoulli density
    $q(\cdot \mid \cdot, \{i,j\})$ is bounded below by
    $\delta^{|V_H|} > 0$ on every binary labeling of $V_H$; in
    particular $q(x_0(P') \mid P, \{i,j\}) > 0$ and
    $q(x_0(P) \mid P', \{i,j\}) > 0$.
\end{enumerate}
\end{lemma}

\begin{proof}
(i) A step that selects pair $\{a,b\}$ relabels vertices only
between districts $a$ and $b$ and changes no other district. If
$P' \ne P$, at least one vertex of $H$ switched between $i$ and $j$,
so districts $i$ and $j$ both changed; any generating pair must
therefore equal $\{i,j\}$. For selectability, note that
$(V_i, V_j)$ and $(V_i', V_j')$ are both bipartitions of the
connected graph $H$ into nonempty parts (all four sets are nonempty
districts of valid plans), and every such bipartition is crossed by
at least one edge of $H$; hence the two districts are adjacent in
both plans.

(ii) The feasible set is $\{x \in [0,1]^{V_H} : \ell \le p^\top x
\le u\}$ with $\ell, u$ given by~\eqref{eq:balance-window}. It
depends only on $(H, \varepsilon, \bar{p})$, not on the endpoint, so
the forward and reverse solves optimize over the same set. It is
nonempty at both endpoints because the window
of~\eqref{eq:balance-window} is precisely the condition that both
sides of the split be valid districts, which $x_0(P)$ and $x_0(P')$
satisfy by validity of $P$ and $P'$. This is where globality is
essential: anchoring the window to the merged-pair half
$(1 \pm \varepsilon)\,p(V_H)/2$ instead lets a valid plan fail its
own pair's window, zeroing a reverse density whose forward
counterpart is positive and destroying reversibility outright. The
objective $\alpha\, x^\top L_H x + \beta\, \|x - x_0\|^2$ has Hessian
$2(\alpha L_H + \beta I) \succ 0$ for $\beta > 0$, so the minimizer
over the (convex, compact, nonempty) feasible set exists and is
unique, and it is a function of $(L_H, x_0, \ell, u)$ only. Since
the kernel $w$ is deterministic, $L_H$ depends only on $H$, and
$x_0$ only on the endpoint and the pair.

(iii) Clipping maps the optimum into $[\delta, 1-\delta]^{V_H}$, so
every Bernoulli factor of the product density lies in
$[\delta, 1-\delta]$ and every outcome in $\{0,1\}^{V_H}$ has
probability at least $\delta^{|V_H|}$.
\end{proof}

\begin{proof}[Proof of Proposition~\ref{prop:qpmh-db}]
Let $K$ denote the transition kernel of Chain E and fix distinct
$P, P' \in \Omega$. If $P$ and $P'$ are not pair-adjacent, no single
step can produce either move: a step changes at most the two selected
districts, so a generating pair would have to be the differing pair;
if the plans differ on more than two districts no pair qualifies, and
if the union $H$ of the two differing districts is disconnected, that
pair is not adjacent in either plan and is never selected. Hence
$K(P, P') = K(P', P) = 0$ and detailed balance holds trivially.

Otherwise let $\{i,j\}$ be the unique generating pair, abbreviate
$q_{\to} = q(x_0(P') \mid P, \{i,j\})$ and
$q_{\gets} = q(x_0(P) \mid P', \{i,j\})$, and note that both are
positive and finite by Lemma~\ref{lem:qpmh-regular}(ii)--(iii).
Because the rounding is single-shot, $q_{\to}$ is the exact
probability that the proposal from $P$ with pair $\{i,j\}$ equals
$P'$: no renormalization by a feasibility event is involved (a
retry-until-feasible loop would replace $q$ by a conditional law
whose state-dependent normalizing constant does not cancel in the
ratio), and clipping \emph{before} sampling makes the sampled law and
the density in~\eqref{eq:qpmh-accept} the same object. Hence
\[
  K(P, P') \;=\; \frac{1}{|\mathrm{bp}(P)|}\; q_{\to}\;
  \alpha(P, P'), \qquad
  \alpha(P, P') = \min\Bigl(1,\;
  \frac{\pi(P')}{\pi(P)}\,
  \frac{q_{\gets}}{q_{\to}}\,
  \frac{|\mathrm{bp}(P)|}{|\mathrm{bp}(P')|}\Bigr),
\]
using $\pi(P')/\pi(P) = e^{-\lambda(|\partial P'| - |\partial P|)}$.
Multiplying through,
\[
  \pi(P)\, K(P, P') \;=\;
  \min\Bigl(
  \frac{\pi(P)\, q_{\to}}{|\mathrm{bp}(P)|},\;
  \frac{\pi(P')\, q_{\gets}}{|\mathrm{bp}(P')|}
  \Bigr),
\]
an expression symmetric under exchanging $(P, q_{\to})$ with
$(P', q_{\gets})$. The identical computation started from $P'$ gives
the same right-hand side, so
$\pi(P) K(P, P') = \pi(P') K(P', P)$.

It remains to check that the rejection pathways do not disturb this
identity. A Bernoulli draw whose split leaves a side empty,
disconnected, or outside the window~\eqref{eq:balance-window} is a
proposal onto a state of zero target mass and is rejected with
probability one, contributing only to the diagonal $K(P, P)$, for
which detailed balance is vacuous; the same applies to in-space
proposals rejected by~\eqref{eq:qpmh-accept}. Rejection on solver
failure is symmetric because, by
Lemma~\ref{lem:qpmh-regular}(ii), the solve is a deterministic
function of $(H, \{i,j\}, \text{endpoint})$, so a failure that
suppresses the move $P \to P'$ suppresses the reverse move
identically. Detailed balance of $K$ with respect to $\pi$ on
$\Omega$ follows, and with it stationarity of
$\pi$~\cite{tierney1998note}.
\end{proof}

\begin{remark}[Scope, verification, and ergodicity]\label{rem:qpmh-scope}
The proof treats the QP solve as exact; in floating point the
minimizer is computed to solver tolerance, and determinism of the
implementation (fixed kernel, no jitter, sorted pair enumeration) is
what the argument actually uses. As an end-to-end check we assembled
the full transition kernel of the implementation on an exhaustively
enumerated $2 \times 3$ grid ($k = 2$, 18 valid states) and verified
$\pi(P)K(P,P') = \pi(P')K(P',P)$ entrywise, with maximum relative
violation $5 \times 10^{-16}$, machine precision. Detailed balance
gives stationarity, not ergodicity: within a selected pair every
binary relabeling of $H$ has positive proposal probability by
Lemma~\ref{lem:qpmh-regular}(iii), but irreducibility of Chain E
across pairs on general $(G, k, \varepsilon)$ is not established, so
convergence statements are for the chain restricted to its
communicating class. On the enumerated $5 \times 5$ fiber the chain
empirically reaches its target mean (\S\ref{sec:exp-longrun}), and
the principled route to global mobility is the $\pi$-Metropolized
mixture described below.
\end{remark}

Chain E trades mobility for exactness: because a proposal must be
feasible in a single Bernoulli draw, moves that flip many units are
rarely accepted, and the chain behaves as a strictly local sampler
(quantified in \S\ref{sec:exp-5x5}). We therefore present it as the
\emph{correctness anchor} of the framework (the member of the family
whose stationarity is provable) rather than as a practical
standalone sampler. The principled route to mobility is precise:
mixing E with a global kernel preserves $\pi$ only if that kernel is
\emph{also} $\pi$-invariant. RevReCom as published is reversible with
respect to its spanning-tree measure, not our tempered $\pi$, so a
naive E/RevReCom mixture would preserve neither; the correct
construction Metropolizes RevReCom's move toward $\pi$ using its
computable Hastings ratio~\cite{cannon2022spanning,autry2023metropolized}
and mixes \emph{that} kernel with E. This is the natural next step
(\S\ref{sec:disc-future}).

\paragraph{On the tempered baseline (Chain A).} Our Tempered ReCom
baseline applies the Metropolis filter
$\min(1, \exp(-\lambda \Delta|\partial P|))$ over the ReCom proposal,
i.e., it \emph{assumes a symmetric proposal}. ReCom's proposal is not
symmetric, the true Hastings ratio involves spanning-tree counts of
the affected districts~\cite{autry2023metropolized,cannon2022spanning}, so Chain A targets $\pi$ only approximately. On the enumerated
$5 \times 5$ fiber this bias is visible and quantifiable
(\S\ref{sec:exp-5x5}: empirical mean cuts $\approx 16.8$ against the
exact target mean $18.56$); it is inherited from the spanning-tree
proposal's own compactness preference, and correcting it requires the
tree-count acceptance ratios of Metropolized forest
recombination~\cite{autry2023metropolized}. We keep the symmetric-$q$
baseline because it is what practitioners run, and we report its bias
rather than hiding it.

\subsection{The integer QP-diffusion (Chain C)}\label{sec:iqp}

The integer counterpart of Definition~\ref{def:qp} replaces
$x \in [0,1]^{|V_H|}$ with $x \in \{0, 1\}^{|V_H|}$ and solves the
resulting mixed-integer convex quadratic program (MIQP) directly.

\paragraph{Definition and solver.}\label{sec:iqp-def}

\begin{definition}[Integer QP-diffusion]\label{def:iqp}
The IQP local refinement of $H$ is
\begin{align}
  x^\star \;=\; \arg\min_{x \in \{0,1\}^{|V_H|}} \;\;
    & \alpha\, x^\top L_H\, x \;+\; \beta\, \|x - x_0\|_2^2 \label{eq:iqp-obj} \\
  \text{subject to} \;\;
    & \ell \;\le\; p^\top x \;\le\; u,
\end{align}
with the same global-ideal balance window $(\ell, u)$ of
Equation~\eqref{eq:balance-window} as the continuous case.
\end{definition}

The objective is convex quadratic and the variable is binary.
Gurobi's MIQP solver handles this directly via branch-and-bound; in
practice, a 5-second time limit per solve is rarely hit on Iowa-sized
merged pairs ($\sim 30$--$60$ precincts).

\paragraph{Subsumption of cut-min; relation to SpecReCom.}\label{sec:iqp-subsumption}

\begin{proposition}[Special cases]\label{prop:iqp-special-cases}
Setting $\alpha = 1$, $\beta = 0$, $w_{ij} \equiv 1$ in
Definition~\ref{def:iqp} gives the classical minimum-cut integer
program subject to the balance window. Hence Chain C subsumes
the discrete cut-minimization MIP as a special case.
\end{proposition}

We are careful \emph{not} to claim the analogous statement for the
continuous case: at $\beta = 0$ the relaxation of
Definition~\ref{def:qp} is degenerate, any feasible constant vector
$c\mathbf{1}$ attains energy zero, so the minimizer carries no cut
information and the sweep is undefined. SpecReCom's Fiedler direction
arises from the Rayleigh quotient \emph{restricted to}
$x \perp \mathbf{1}$, a constraint our proximal program does not
impose; it is the fidelity term ($\beta > 0$) that breaks the constant
degeneracy here, with $x_0$ playing the symmetry-breaking role that
orthogonality plays in spectral partitioning. The precise relationship
is: as $\beta \to 0^+$ with $x_0$ non-constant, the minimizer's
projection off $\mathbf{1}$ aligns with the low-frequency eigenspace
of $L_H$ (Equation~\eqref{eq:freq-response} with
$\beta/(\alpha\lambda_k + \beta) \to 0$ fastest at high frequencies),
so Chain B approaches a Fiedler-like split \emph{in direction} while
retaining balance feasibility, a BalSpecReCom-style variant rather
than SpecReCom itself. Chain D implements that variant directly.

\paragraph{Locality via the fidelity term.}\label{sec:iqp-locality}

For $\beta = 0$, the IQP optimum is a minimum cut of $H$ subject only
to population balance and may differ arbitrarily from $x_0$. For
$\beta > 0$, the optimum is biased toward $x_0$: any flip from $x_0$
incurs a $\beta$-penalty per bit changed. This has three coupled
benefits over $\beta = 0$:

\begin{enumerate}
  \item \textbf{Locality.} The proposed bipartition stays geographically
    close to the current one, so the chain explores neighboring
    plans rather than randomly resetting on each merge.
  \item \textbf{Contiguity-friendliness.} Because $x_0$ encodes a
    contiguous bipartition and the IQP optimum is biased toward $x_0$,
    the optimum is mostly contiguous already; orphan-flip repair
    needs only to fix a small number of boundary-flipping vertices.
  \item \textbf{Lower rejection.} Post-rounding rejection on Iowa
    drops from 5--15\% ($\beta = 0$) to $<3\%$ ($\beta = 1$) with
    repair, while cut quality is preserved.
\end{enumerate}

\paragraph{One step of Chain C.}\label{sec:chain-c-step}
A single step of Chain C wraps an outer Tempered ReCom move
(Algorithm~\ref{alg:recom} plus the Metropolis--Hastings filter of
\S\ref{sec:hybrid-mh}) around $n_{\mathrm{iqp}}$ deterministic
IQP local-refinement passes. Each inner pass picks an adjacent district
pair, solves the IQP on the merged subgraph, and applies orphan-flip
repair; if the repair fails the inner pass leaves the plan unchanged.
The full step is given as Algorithm~\ref{alg:chain-c}.

\begin{algorithm}[H]
\caption{One step of Chain C}\label{alg:chain-c}
\KwIn{Current plan $P$, parameters $\alpha, \beta, n_{\mathrm{iqp}}, \lambda, \varepsilon$.}
\KwOut{Next plan $P'$.}
$P_0 \gets $ Tempered ReCom step (Chain A) applied to $P$\;
$P_{\mathrm{curr}} \gets P_0$\;
\For{$t = 1$ \KwTo $n_{\mathrm{iqp}}$}{
  Select adjacent district pair $(V_i, V_j)$ uniformly from boundary pairs of $P_{\mathrm{curr}}$\;
  Form $H = G[V_i \cup V_j]$, build kernel-weighted $L_H$\;
  Solve IQP (Definition~\ref{def:iqp}) for $x^\star$\;
  Apply orphan-flip repair (Algorithm~\ref{alg:orphan-flip})\;
  \uIf{repaired bipartition feasible and connected}{
    Update $V_i, V_j$ in $P_{\mathrm{curr}}$ to the new bipartition\;
  }\Else{
    $P_{\mathrm{curr}}$ unchanged\;
  }
}
\Return $P' \gets P_{\mathrm{curr}}$.
\end{algorithm}

\paragraph{What Chain C is, and is not.}\label{sec:chain-c-honest}
The IQP loop is \emph{deterministic} given the random pair choices, there is no Metropolis acceptance inside the loop, and the outer
acceptance applies only to the Tempered ReCom move. We are explicit
about the consequence: \textbf{Chain C is not a sampler for the
tempered target} $\pi(P) \propto \exp(-\lambda|\partial P|)$, and no
completion of the construction can make it one. A deterministic map
$T$ composed after a $\pi$-invariant kernel preserves $\pi$ only if
$T$ is itself $\pi$-preserving (essentially a $\pi$-preserving
involution~\cite{tierney1998note}) whereas the IQP refinement is a
many-to-one, cut-decreasing contraction onto local minima: the
opposite of an involution. Nor can the refinement be folded into the
proposal and Metropolized, since the reverse move would require the
argmin at $P'$ to output the strictly worse plan $P$, an event of
probability zero. The distributional signature of this bias is visible
in our own ground-truth data (\S\ref{sec:exp-5x5}): Chain C's cut
marginal is shifted and sharply concentrated relative to Chain A's,
despite the identical outer kernel.

We therefore present Chain C as what it is: a \emph{matheuristic
ensemble generator} in the lineage of ReCom-based
optimization~\cite{cannon2023voting,brous2025recom}, an iterated
local search whose stochastic component (the tempered outer move)
maintains diversity while the embedded MIQP drives compactness. This
is a legitimate and, as \S\ref{sec:experiments} shows, empirically
powerful tool for producing high-quality compact plans and for
seeding or stress-testing ensembles. What it does not provide is a
defensible null distribution for outlier analysis; for that use case
the field's requirements are exactly the known-target
property~\cite{cannon2022spanning,mccartan2023sequential} that
Chain E (\S\ref{sec:qp-mh}) preserves and Chain C deliberately trades
away for optimization power.

\paragraph{Pareto trade-off: continuous vs.\ integer.}\label{sec:iqp-pareto}

\begin{table}[H]
\centering
\small
\begin{tabular}{lll}
\toprule
 & Chain B (continuous QP) & Chain C (integer IQP) \\
\midrule
Variable & $x \in [0,1]^n$ & $x \in \{0,1\}^n$ \\
Solver & Convex QP (fast) & Convex MIQP (branch-and-bound) \\
Per-step cost & $O(\text{nnz}(L_H))$ & $O$(MIQP, $\le 5$s limit) \\
Cut quality (Iowa) & $-4.0\%$ vs A (n.s.) & $-21.1\%$ vs A ($d = -2.65$) \\
Contiguity & after rounding + repair & after repair (mostly already integer) \\
\bottomrule
\end{tabular}
\caption{Chain B and C trade-off (Table~\ref{tab:iowa-headline}
  numbers). B is $\sim 40\times$ faster per step; only C moves the
  compactness needle on Iowa. Both kernel-flexible.}\label{tab:b-vs-c}
\end{table}

%════════════════════════════════════════════════════════════════════════════════
\section{Experiments}\label{sec:experiments}
%════════════════════════════════════════════════════════════════════════════════

We evaluate the six chains of Table~\ref{tab:five-chains} on three
datasets: the $5 \times 5$ grid (full enumeration, ground-truth
target), the Iowa county dual graph (real geography, 99 counties,
$k = 4$, $\varepsilon = 0.05$), and \draft{NC precincts (k=4
sub-state; pending Sprint 3 completion)}. All hyperparameters are
locked in Appendix~\ref{app:hyperparams}; raw cut traces, per-step
balance audits, and configuration stamps are released in the project
repository.

\paragraph{Statistical protocol.} Chains are autocorrelated, so we do
not report i.i.d.\ test statistics on raw traces. For every
comparison we report the mean difference, Cohen's $d$ on the pooled
post-burn-in standard deviation, and a $z$-test whose standard errors
use effective sample sizes from Geyer's initial-positive-sequence
estimator of the integrated autocorrelation time; ESS values (tens,
not hundreds, per 400--900 retained steps) are reported alongside.
Each chain runs on its own seeded RNG stream. The tables below are
from a single seed (42); \draft{a 10-seed replication with
between-seed uncertainty is queued and will replace the single-seed
columns before submission.}

\paragraph{Balance audit.}\label{sec:exp-balance} Every experiment
records, at every step, the maximum relative deviation of any district
population from the global ideal. All six chains satisfy the
$\varepsilon = 0.05$ bound at every step of every run reported here
(Iowa maxima 0.0498--0.0500). We flag this because it is not automatic:
GerryChain's raw \texttt{recom} proposal validates only one side of
each tree cut, and an earlier version of our own local steps anchored
balance to the merged-pair half rather than the global ideal
(\S\ref{sec:qp-model}); either slip lets district populations drift
to $\sim$4$\times$ the nominal tolerance on Iowa before the audit
catches it. We recommend per-step balance auditing as standard
practice for any custom chain harness that bypasses GerryChain's
validator pipeline.

\subsection{Multi-kernel showcase on Iowa}\label{sec:exp-multikernel}

We run Chain B and Chain C on Iowa with three kernels and visualize
the resulting partitions. The three kernels are uniform
($k_{\mathrm{unif}}$), perimeter ($k_{\mathrm{perim}}$,
$w_{ij}$ = shared-boundary length from \texttt{shared\_perim}), and
density ($k_{\mathrm{dens}}$, $w_{ij} = |\partial v_i \cap \partial v_j| \cdot \exp(-(\rho_i - \rho_j)^2 / \sigma^2)$
with $\rho_v = \mathrm{TOTPOP}/\mathrm{area}$ normalized to unit
mean and $\sigma = 1$).

\begin{figure}[H]
  \centering
  \IfFileExists{fig/multikernel_iowa.png}{%
    \includegraphics[width=\textwidth]{fig/multikernel_iowa.png}%
  }{%
    \fbox{\begin{minipage}{0.9\textwidth}\vspace{2cm}\centering
      \ttfamily\small\textcolor{orange}{[Missing figure:
      \texttt{fig/multikernel\_iowa.png}]}\vspace{2cm}\end{minipage}}%
  }
  \caption{Final Chain C plans on Iowa under the three kernels
    (seed 42, 500 steps, identical constraints; only the Laplacian
    edge weights differ). Panel annotations give the unweighted cut
    count and the geographic boundary length (sum of shared-perimeter
    weights over cut edges). The plans differ in 24 to 30 of the 99
    counties under optimal district relabeling: the kernel directly
    encodes geographic signal at the proposal level.}
  \label{fig:multikernel-iowa}
\end{figure}

\begin{table}[H]
\centering
\small
\begin{tabular}{lrrrr}
\toprule
Kernel & B mean cuts & C mean cuts & C plan boundary
  & C cut-edge $|\Delta\rho|$ ratio \\
\midrule
$k_{\mathrm{unif}}$  & 38.15 & 31.39 & 10.24 & 1.00 \\
$k_{\mathrm{perim}}$ & 39.00 & 33.39 & 9.82  & 1.51 \\
$k_{\mathrm{dens}}$  & 38.09 & 36.84 & 9.35  & 2.41 \\
\bottomrule
\end{tabular}
\caption{Three kernels on Iowa ($k = 4$, $\varepsilon = 0.05$,
  500 steps, burn-in 100, seed 42, corrected balance constraint).
  Mean cuts are post-burn-in averages of the unweighted cut count;
  ``C plan boundary'' is the geographic boundary length of the final
  Chain C plan (sum of \texttt{shared\_perim} over cut edges, in the
  dual graph's units); the last column is the mean population-density
  gap across the final plan's cut edges divided by the graph-wide
  mean gap, so $1.00$ means density-blind boundaries and larger
  values mean boundaries concentrated on urban--rural transitions.
  Under non-uniform kernels the unweighted cut count is not the
  objective the kernel optimizes, so the cut columns must be read
  together with the kernel-specific columns.}\label{tab:multikernel-iowa}
\end{table}

\textbf{Positioning of the kernel showcase.} SpecReCom uses randomized
edge weights only to break spectral ties; region-aware ReCom and
county-weighted trees accept indicator surcharges; none accepts an
arbitrary user-supplied kernel with a locality anchor. The claim this
experiment supports is deliberately scoped: kernels as first-class
proposal inputs, with characterized behavior, not priority over
all weighted proposals (\S\ref{sec:related-weighted}).

\textbf{What each kernel does to the plan.} The three final Chain C
plans in Figure~\ref{fig:multikernel-iowa} are genuinely different
districtings, not perturbations of one another: under the best
district relabeling they disagree on 30 of 99 counties (uniform vs.\
perimeter), 25 (uniform vs.\ density), and 24 (perimeter vs.\
density). Each kernel improves exactly the quantity it encodes. The
uniform kernel minimizes the unweighted cut count and achieves the
fewest cut edges (31) but the longest geographic boundary. The
perimeter kernel prices each edge by its shared boundary length, and
its final plan trades one extra cut edge for a $4.1\%$ shorter
geographic boundary. The density kernel prices edges by shared
boundary length damped across density gaps, and its plan shortens the
geographic boundary by $8.7\%$ while concentrating cuts on
urban--rural transitions: the mean population-density gap across its
cut edges is $2.41\times$ the graph-wide mean, versus $1.00\times$
(exactly density-blind) for the uniform kernel. In map terms, the
density plan's largest-gap cut edges are precisely the borders of
Polk (Des Moines), Linn (Cedar Rapids), and Johnson (Iowa City)
counties, the state's three main urban cores, rather than edges
through uniformly rural blocks.

\textbf{Reading the cut-count column honestly.} Under non-uniform
kernels the unweighted cut count worsens, from 31.39 (uniform) to
33.39 (perimeter) and 36.84 (density) for Chain C, and it should:
raw cut count is no longer the objective. The weighted metrics
improve in lockstep, which is the point of the showcase. Chain B is
nearly kernel-insensitive (38.1--39.0 mean cuts across all three),
consistent with its Iowa null result in
Table~\ref{tab:iowa-headline}: the continuous relaxation plus
rounding does not transmit the kernel's signal into the final integer
plan the way the exact integer solve does.

These are single-seed runs (the maps are final-plan snapshots, the
table means are post-burn-in averages over 500 steps), so we read
them as a demonstration of direction, not as calibrated effect sizes;
the direction is unambiguous because each kernel's improvement shows
up in its own metric on the same seed and the same constraint set.

\subsection{Compactness comparison across datasets}\label{sec:exp-compactness}

\paragraph{Iowa (6 chains).}\label{sec:exp-iowa}

\begin{table}[H]
\centering
\small
\begin{tabular}{lrrllr}
\toprule
Chain & Mean cuts & Std & $\Delta$ vs A (Cohen's $d$) & ESS-$z$ $p$ & Wall-clock \\
\midrule
V: Vanilla ReCom         & 41.39 & 4.78 & ---                & ---     & 1.2\,s \\
A: Tempered ReCom        & 39.76 & 4.32 & ---                & ---     & 1.1\,s \\
B: QP-diffusion mix      & 38.15 & 4.28 & $-4.0\%$ ($-0.37$) & $0.24$  & 2.7\,s \\
\textbf{C: IQP local ($n_{\mathrm{iqp}} = 10$)}
                         & \textbf{31.39} & \textbf{1.15}
                         & \textbf{$-21.1\%$ ($-2.65$)} & $< 10^{-12}$ & 92.6\,s \\
D: Spectral mix          & 39.51 & 4.70 & $-0.6\%$ ($-0.06$) & $0.85$  & 1.5\,s \\
E: QP-MH exact           & 55.21$^{\dagger}$ & 7.88 & n/a (not converged) & n/a & 19.8\,s \\
\bottomrule
\end{tabular}
\caption{Iowa cut counts ($k = 4$, $\varepsilon = 0.05$, 500 steps,
  burn-in 100, seed 42, per-step balance audit passed by all chains).
  Chain C delivers a 21.1\% reduction in mean cuts versus Tempered
  ReCom with the tightest distribution; Chains B and D are
  statistically indistinguishable from A once balance is enforced
  correctly and autocorrelation is accounted for. $^{\dagger}$Chain E
  is not a competitor on compactness: it is the exact sampler for the
  tempered target, its trajectory is still rising at 500 steps
  (ESS $\approx 3$), so we report its mean as a \emph{non-equilibrium
  lower bound} on the target's cut level, with no test statistic.
  Even as a lower bound it shows the target
  $\pi \propto \exp(-0.1\,|\partial P|)$ is far \emph{less} compact
  than the output of any ReCom-family chain.}\label{tab:iowa-headline}
\end{table}

Three honest readings of Table~\ref{tab:iowa-headline}. First, the
integer refinement is the only local kernel whose Iowa gain survives
rigorous accounting: $-21.1\%$ with $d = -2.65$ is an enormous effect
by any convention, and it is robust to the balance-constraint
correction that erased Chain B's previously-measured advantage (an
earlier draft reported a significant $-3.2\%$ for B under the loose
merged-pair balance anchor; with the global window enforced, B's
point estimate swings between $-0.3\%$ and $-4.0\%$ across equivalent
reruns and never approaches significance at these effective sample
sizes, a cautionary tale about how constraint bugs can masquerade as
method gains).
Second, Chain D's parity with A and V supports our claim that the
spectral cut alone, without the fidelity term, does not improve
ReCom-family compactness on heterogeneous county graphs. Third, Chain
E's trajectory exposes a fact the field's tempered-chain practice
obscures: taking E's non-equilibrium mean as a lower bound, every
ReCom-based chain here sits at least $\sim$15--25 cuts \emph{below}
the tempered target's own mean, because the spanning-tree proposal is
itself strongly compactness-biased. ``Tempered ReCom targets $\pi$''
is, on this evidence, a statement about the filter, not the samples.

\begin{figure}[H]
  \centering
  \IfFileExists{../diffusion_model/recom/iowa_hybrid_test_out/target_matched_iowa.png}{%
    \includegraphics[width=\textwidth]{../diffusion_model/recom/iowa_hybrid_test_out/target_matched_iowa.png}%
  }{%
    \needdata{include iowa\_hybrid\_test\_out/target\_matched\_iowa.png}%
  }
  \caption{Cut-edge trace and post-burn-in histogram for the six
    chains on Iowa. Chain C concentrates tightly at low cut counts;
    Chain E drifts upward toward the tempered target's own cut
    level.}
  \label{fig:iowa-trace}
\end{figure}

The wall-clock column completes the practitioner picture: Chain C's
compactness costs $\sim 80\times$ Vanilla ReCom per step at
$n_{\mathrm{iqp}} = 10$ (about 0.2\,s per step on Iowa), which is
acceptable for ensembles of thousands of plans but material at
precinct scale. Chain C's tight standard deviation (1.15 vs.\ 4.08
for A) is the signature of an optimizer concentrating on IQP local
minima rather than a sampler's honest spread, consistent with its
matheuristic reading (\S\ref{sec:chain-c-honest}).

\paragraph{5$\times$5 ground truth: coverage and distributional bias.}\label{sec:exp-5x5}

The $5 \times 5$ partition fiber ($k = 5$, exact balance, 4\,006
partitions; \S\ref{sec:fiber-corridors}) admits full enumeration, so
every distributional quantity is exactly computable. Two disclaimers
shape the design of this experiment. First, \emph{uniform-on-fiber is
not the target of any chain in this paper}; we report TV-to-uniform
only as a descriptive concentration measure. Second, with $\sim$900
retained correlated draws over 4\,006 states, even a perfect
i.i.d.\ uniform sampler would show TV $\gtrsim 0.78$, short-run TV
values mostly restate the coverage column. The informative
comparisons are (i) the empirical cut \emph{mean} against the exactly
computable target mean, which converges at rate
$\mathrm{ESS}^{-1/2}$ and does not suffer the support problem, and
(ii) long runs (\S\ref{sec:exp-longrun}).

\begin{table}[H]
\centering
\small
\begin{tabular}{lrrrr}
\toprule
Chain & Mean cuts & Std & Unique / 4006 & TV(emp, uniform) \\
\midrule
V: Vanilla ReCom     & 16.91 & 1.06 & 341 / 4006 (8.5\%) & 0.91 \\
A: Tempered ReCom    & 16.75 & 0.95 & 297 / 4006 (7.4\%) & 0.93 \\
B: QP-diffusion mix  & 16.62 & 0.88 & 164 / 4006 (4.1\%) & 0.96 \\
C: IQP local         & 16.15 & 0.52 & 151 / 4006 (3.8\%) & 0.96 \\
D: Spectral mix      & 16.58 & 0.82 & 253 / 4006 (6.3\%) & 0.94 \\
E: QP-MH exact       & 16.27 & 0.44 & \phantom{00}6 / 4006 (0.1\%) & 1.00 \\
\bottomrule
\end{tabular}
\caption{$5 \times 5$ short-run comparison ($k = 5$, exact balance,
  1\,000 steps, burn-in 100, seed 42). The exact target mean is
  $\E_\pi[|\partial P|] = 18.56$ at $\lambda = 0.10$: \emph{every}
  chain sits far below it at this horizon, including the exactly
  reversible Chain E, which at 1\,000 steps has visited only 6 states, exactness of the kernel does not buy mixing speed. Compactness
  gains trade against coverage: C covers 3.8\% of the fiber against
  V's 8.5\%.}\label{tab:5x5-headline}
\end{table}

\textbf{Cost of compactness, honestly stated.} Chains B and C
sacrifice roughly half of V's fiber coverage at this horizon. The
gap between every chain's empirical mean and the target mean 18.56
quantifies the compactness bias \emph{shared} by the whole
ReCom-derived family, and on the enumerated fiber we can identify its
origin exactly. Using the Matrix--Tree theorem we compute, for each
of the 4\,006 partitions, the spanning-tree-count weight
$\prod_i \tau(V_i)$, the measure that the RevReCom analysis
identifies as the natural target of tree-based
recombination~\cite{cannon2022spanning,autry2023metropolized}. The
match is striking: the tree-weighted measure has
$\E[|\partial P|] = 16.88$ against Vanilla ReCom's empirical 16.91,
and the tree-weighted-and-tempered measure
$\prod_i \tau(V_i)\, e^{-\lambda|\partial P|}$ has
$\E[|\partial P|] = 16.78$ against Chain A's empirical 16.75. In
other words, the ``tempered gap'' is not a mystery: on this test bed
the ReCom-family chains sit almost exactly on their
tree-count-weighted quasi-targets, 1.7--1.8 cuts below the nominal
tempered distribution, quantitatively confirming the
spanning-tree-measure theory on fully enumerated ground truth. Our
refinements deepen that pre-existing bias only modestly (Chain C
lands 2.4 cuts below the nominal target). The much larger relative
gain on Iowa ($-21.1\%$ vs.\ $-3.6\%$ here) reflects the room that
heterogeneous county populations leave for optimization; the uniform
grid leaves little, so we read the $5 \times 5$ results as
calibration, not as a demonstration of power.

\paragraph{Long-run validation.}\label{sec:exp-longrun}
We run Chain A for $2 \times 10^5$ steps and Chain E for $10^5$ steps
on the enumerated fiber (seed 42, 10\% burn-in) and compare empirical
cut means against the exact target mean
$\E_\pi[|\partial P|] = 18.557$. The two chains test two predictions,
and both are confirmed.

\emph{Chain A: the symmetric-$q$ bias persists asymptotically.} At
$2 \times 10^5$ steps the tempered chain has visited $3\,524/4\,006$
states, and its empirical mean is $16.73$: the $1.8$-cut gap from the
nominal target does not shrink with run length. It is not drift; it is
convergence to the wrong (tree-weighted) distribution: the
Matrix--Tree prediction for the tree-weighted-and-tempered measure is
$16.78$, and the chain lands within $0.05$ cuts of it. The
symmetric-$q$ filter converges, quantitatively, to the
spanning-tree-weighted analogue of its target, exactly as the
RevReCom analysis predicts~\cite{cannon2022spanning}.

\emph{Chain E: the exact sampler converges in mean to its target.} At
$10^5$ steps Chain E's empirical mean is $18.39$, within $0.17$ cuts
(under $1\%$) of the exact target mean, with the gap shrinking from
$2.3$ cuts at the $10^3$-step horizon. This is end-to-end empirical
validation of the detailed-balance construction of \S\ref{sec:qp-mh}
on a nontrivial state space, complementing the machine-precision
kernel check of Remark~\ref{rem:qpmh-scope}. Convergence of the full
distribution is far slower, as expected for a sticky local chain: E
has visited only $403/4\,006$ states and its TV distance to $\pi$
remains $0.90$, so the mean (which converges at the ESS rate on a
bounded function) is the honest headline statistic, not TV.

\subsection{Ablations and sensitivity}\label{sec:exp-ablations}

\draft{Caveat: the ablation tables in this subsection were produced
with the earlier merged-pair balance anchor
(\S\ref{sec:exp-balance}) and are retained as indicative pending
regeneration with the corrected constraint; the headline tables above
are from the corrected code.}

\paragraph{Diversity controls: Jitter and Tabu.}\label{sec:exp-diversity}

\begin{table}[H]
\centering
\small
\begin{tabular}{lrrrr}
\toprule
Configuration & Mean cuts & Std & Coverage & Wall \\
\midrule
Baseline ($n_{\mathrm{iqp}} = 5$)               & 31.88 & 2.28 & 4.5\% & 39.7\,s \\
+ Jitter ($\mathcal{U}[1,2]$ on edge weights)   & 31.76 & 1.77 & 5.7\% & 45.2\,s \\
+ Tabu (size 3)                                  & \textbf{31.29} & \textbf{1.26} & 4.6\% & \textbf{37.6\,s} \\
\bottomrule
\end{tabular}
\caption{Diversity controls on Iowa ($n_{\mathrm{iqp}} = 5$). Tabu
  delivers the best combination of cut quality, distribution
  tightness, and runtime.}\label{tab:diversity-iowa}
\end{table}

Tabu, in particular, tightens the cut-count distribution
(std $2.28 \to 1.26$, a 44\% reduction) without sacrificing coverage
and at a modest runtime saving. Both controls are configurable knobs
for the practitioner trading compactness against representativeness.

\note{Add a paragraph explaining \emph{why} Jitter and Tabu help:
Jitter breaks ties between equally-good cuts, encouraging the chain
to explore equivalent configurations rather than locking onto one
deterministic optimum (which would happen with $\beta > 0$ and a tight
MIQP); Tabu forbids the same merged-pair from being picked
consecutively, forcing the chain to vary which boundary it refines.
Frame them as orthogonal: Jitter affects \emph{which cut} on a given
pair, Tabu affects \emph{which pair} the chain picks next.}

\paragraph{Parameter sensitivity: $n_{\mathrm{iqp}}$, $\alpha$, $\beta$.}\label{sec:exp-sensitivity}

\begin{table}[H]
\centering
\small
\begin{tabular}{lrrr}
\toprule
Configuration & Mean cuts & $\Delta$ vs A & Wall-clock \\
\midrule
$n = 3, \alpha = 10, \beta = 1$  & 32.37 & $-19.6\%$ & 24\,s \\
$n = 5, \alpha = 10, \beta = 1$  & 31.88 & $-20.8\%$ & 40\,s \\
\textbf{$n = 10, \alpha = 10, \beta = 1$} & \textbf{30.66} & \textbf{$-23.9\%$} & 103\,s \\
$n = 3, \alpha = 50, \beta = 1$  & 32.38 & $-19.6\%$ & 32\,s \\
$n = 3, \alpha = 10, \beta = 0$  & 32.31 & $-19.8\%$ & 31\,s \\
$n = 5, \alpha = 10, \beta = 0$  & 31.45 & $-21.9\%$ & 46\,s \\
\bottomrule
\end{tabular}
\caption{Chain C parameter sensitivity on Iowa. $n_{\mathrm{iqp}}$ is
  the dominant lever; $\alpha$ saturates; $\beta$ has minor effect
  on cut quality but governs the locality structure exploited by
  orphan-flip repair.}\label{tab:sensitivity}
\end{table}

\note{Add 1--2 paragraphs interpreting Table~\ref{tab:sensitivity}:
(1) Why does $n_{\mathrm{iqp}}$ scale roughly linearly with runtime
but \emph{sub-linearly} with cut reduction (i.e., diminishing returns)?
(2) Why does $\alpha$ saturate at 10? At what $\alpha$ does
oversmoothing begin?
(3) Why is $\beta = 0$ \emph{nearly as good} on cuts but with no
contiguity guarantee, giving up locality buys almost nothing on
this metric but costs the rejection-rate advantage.}

\subsection{Pending experiments}\label{sec:exp-pending}

\paragraph{North Carolina precinct validation.}\label{sec:exp-nc}
\needdata{Pending Sprint 3 (PAPER\_PLAN \S7.5). NC precinct (2\,692
nodes, $k = 4$). Expect Chain C per-step cost $\sim 5$--30\,s; total
$\sim 10$--30 min for 100 steps. Cross-state replication of the
Iowa compactness gain.}

%════════════════════════════════════════════════════════════════════════════════
\section{Discussion}\label{sec:discussion}
%════════════════════════════════════════════════════════════════════════════════

\subsection{What each member of the family is for}\label{sec:disc-roles}

The three new chains occupy deliberately different points on the
correctness--power spectrum, and the honest summary of our results is
that no single point dominates. \textbf{Chain C} is where the
practical power lives: a $21.1\%$ cut reduction with Cohen's
$d = -2.65$ on Iowa is the kind of improvement practitioners can see
on a map, and the MIQP's optimality within each merged pair is what
makes it robust where the continuous relaxation (Chain B) fails. But
that power is an optimizer's power. Its tight output spread
(std $1.15$ vs.\ $4.08$) is concentration on IQP local minima, and we
have shown (\S\ref{sec:chain-c-honest}) that no lemma can reconcile
its deterministic contraction with a stationarity claim. Use Chain C
to produce high-quality compact plans, to seed other chains, or to
stress-test the compactness frontier of a state; do not use it to
build a null distribution for outlier analysis. \textbf{Chain E} is the
opposite pole: provably exact for its target and, by the same token,
brutally honest about what that target is. Its non-equilibrium Iowa
trajectory climbing past $55$ cuts (a lower bound on the target's
level), with every other chain sitting at least $15$--$25$ cuts
below, is to our knowledge the first direct empirical display of how
far the ReCom family's outputs sit from the tempered distribution its
Metropolis filters nominally target; the enumerated fiber, where the
gap is exactly computable, confirms the direction of the effect. \textbf{Chain
B} occupies the middle and, after correcting the balance-constraint
bug that had inflated its measured advantage, we report it as a
negative result on Iowa: the continuous relaxation plus sweep
rounding adds nothing significant over Tempered ReCom on
heterogeneous county graphs. We keep it in the paper because the
QP/IQP contrast is informative (integrality, not spectral
smoothing, is doing the work) and because its per-step cost is
$40\times$ lower than C's.

Why does the fidelity term matter? The evidence is the Iowa
compactness contrast, not the jump-rate table. Chain D, our
balance-constrained variant of the SpecReCom idea with no $\beta$
term, is statistically indistinguishable from Vanilla and Tempered
ReCom on Iowa (Table~\ref{tab:iowa-headline}): the spectral cut alone
re-splits each merged pair from scratch and discards the structure the
chain has already found, whereas the anchored integer refinement
improves on it decisively. The community step-trace diagnostic
(\S\ref{sec:fiber-corridors}) makes a complementary, weaker point
about locality in general: pure tree-move chains cross community
boundaries at 14--16\% per step, near the community-blind baseline,
while every mixture with an optimization-guided local move (anchored
or not) drops the rate to 6--9\%. The hybrid design, ReCom for
across-community mobility and anchored refinement for
within-community descent, follows from the two observations
together.

\subsection{On the macro-structure findings}\label{sec:disc-throats}

The corrected picture of \S\ref{sec:fiber-corridors} has two
consequences worth stating plainly. First, it relocates the blame for
slow local mixing. If plan space were islands joined by narrow
corridors, slow traversal would be a topological inevitability and no
local proposal could fix it; since the fibers are in fact well
connected (a fifth to a quarter of edges cross communities, and every
pair of large communities shares a wide interface), slow traversal is
a property of the dynamics, and better local proposals are not
hopeless. Chain E crossing a community boundary 3 times in 1\,000
steps despite 18\% of edges being available is the extreme datum: its
bottleneck is acceptance of large moves, not the graph. Second, the
finding sharpens what a mixing obstruction for these state spaces
would have to look like. Worst-case hardness
results~\cite{akitaya2022reconfiguration,charikar2023complexity}
permit bottlenecked instances, and our grids may simply be too small
or too symmetric to exhibit them; enumerating a fiber with a genuine
conductance bottleneck (or proving none exists for grid families)
would be a valuable follow-up, and the anti-bottleneck statistics of
Table~\ref{tab:narrow-throats} give the template for checking.

\subsection{Limitations}\label{sec:disc-limits}

\textbf{(L1) Chain B adds nothing significant on Iowa.} The
continuous QP-diffusion's earlier measured gain was an artifact of a
balance-constraint bug; corrected, its estimate is not
distinguishable from zero ($-4.0\%$, $p = 0.24$, with point estimates
swinging between $-0.3\%$ and $-4.0\%$ across equivalent reruns at
ESS $\approx 30$). The relaxation appears too loose on graphs with
highly variable vertex populations.

\textbf{(L2) No theoretical mixing-time bound.} Like every existing
redistricting chain, we lack a polynomial mixing-time guarantee, and
Chain E demonstrates that even an exactly reversible member of the
family can be an extremely slow mixer. We rely on empirical
diagnostics, ESS, the community step-trace diagnostic, and enumerated
ground truth.

\textbf{(L3) Chain A's symmetric-$q$ approximation is not a small
bias.} On the enumerated fiber the tempered chain's empirical mean
($16.75$) sits $1.8$ cuts below the exact target mean ($18.56$), and
Chain E's Iowa trajectory suggests the gap is larger, not smaller, on
real graphs. Correcting it requires spanning-tree-count acceptance
ratios~\cite{autry2023metropolized} or a RevReCom outer
kernel~\cite{cannon2022spanning}.

\textbf{(L4) Chain E mixes too slowly to use alone.} Exactness comes
from single-shot rounding, which caps move sizes; the chain visited
6 of 4\,006 fiber states in 1\,000 steps. It validates the corrected
proposal and calibrates other chains' bias; it is not a production
sampler.

\textbf{(L5) Integer QP is computationally expensive.} Chain C costs
$\sim 100\times$ Vanilla ReCom per step at $n_{\mathrm{iqp}} = 10$ on
Iowa. For precinct-scale graphs, time limits with warm starts (and
possibly coarsened merged-pair subproblems) will be necessary; the
NC experiment (\S\ref{sec:exp-nc}) is the pending scalability test.

\textbf{(L6) Single-seed tables.} All reported tables are seed-42
runs with per-chain RNG streams; the 10-seed replication is queued.
Effect sizes of Chain C's magnitude ($d < -2.5$) will not be
overturned by seed variation, but the smaller contrasts (B vs.\ A,
D vs.\ A) should be read as provisional.

\subsection{Future directions}\label{sec:disc-future}

\begin{enumerate}
  \item \textbf{A reversible hybrid.} Mix Chain E's exact local
    kernel with a global move that is Metropolized toward the
    \emph{same} target $\pi$ (e.g., a RevReCom-style tree move
    accepted with its computable spanning-tree-count Hastings
    ratio~\cite{cannon2022spanning,autry2023metropolized}) giving a
    chain that is mobile across communities and exactly stationary for a
    stated target. (A naive mixture of kernels reversible for
    \emph{different} targets preserves neither.) A proof-of-concept on
    the enumerated $5 \times 5$ fiber, where convergence is checkable
    exactly, is the single most promising follow-up.
  \item Replica exchange or parallel tempering between Chain C and
    sampling chains, in the spirit of multiscale parallel
    tempering~\cite{chuang2025multiscale}, to combine compactness with
    coverage.
  \item Integration with the cycle walk~\cite{deford2025cycle} or the
    balanced up-down walk~\cite{akitaya2026bud} as outer kernels.
  \item Benchmarking Chain C against exact cut-minimization optima on
    the public instances of Validi and
    Buchanan~\cite{validi2022political}, positioning it properly in
    the matheuristic literature~\cite{brous2025recom}.
  \item Multi-scale diffusion kernels (heat kernel, commute-time
    kernel), and a kernel-faithful Cheeger sweep that scores
    candidates by weighted cut (\S\ref{sec:rounding}).
\end{enumerate}

%════════════════════════════════════════════════════════════════════════════════
\section{Conclusion}\label{sec:conclusion}
%════════════════════════════════════════════════════════════════════════════════

This paper introduced a kernel-flexible variational framework for
local refinement in redistricting Markov chains: one Laplacian energy
$\alpha\,x^\top L x + \beta\,\|x - x_0\|^2$ under global population
balance, three solution regimes. The continuous relaxation with sweep
rounding (Chain B) is fast but adds no significant compactness on
heterogeneous graphs; the integer program (Chain C) is a matheuristic
that delivers a $21.1\%$ cut reduction (Cohen's $d = -2.65$) over
Tempered ReCom on Iowa while deliberately forgoing any stationarity
claim; and the Metropolis-corrected randomized-rounding variant
(Chain E) is, to our knowledge, the first optimization-based proposal
in redistricting MCMC with an exact detailed-balance argument. A
structural diagnostic of three fully enumerated partition fibers
finds genuine community structure but no narrow corridors, so global
tree moves and anchored local refinements are complementary for
dynamical rather than topological reasons,
and per-step balance and ESS audits keep the empirical claims honest,
including against ourselves: two constraint bugs we found and fixed
during this work had inflated earlier measurements, and Chain E
reveals that the whole ReCom family samples far more compactly than
its nominal tempered targets.

For the field, we hope two ideas outlast the specific chains. First,
kernel choice (shared perimeter, density gradients, administrative
boundaries) can be a first-class modeling decision \emph{inside}
the proposal rather than a post-hoc filter, provided the induced
distribution is either corrected (Chain E) or the method is honestly
labeled an optimizer (Chain C). Second, enumerated fibers are cheap,
severe test beds: a chain's distance from its own computable target
on 4\,006 states is a more informative diagnostic than any p-value on
a single trace. The most promising next step is the reversible
hybrid: Chain E's exact local kernel mixed with a global tree move
Metropolized toward the same target via its computable
spanning-tree-count ratio, yielding a mobile chain whose stationary
distribution is known, the property that court-facing ensemble
analysis ultimately requires.

%════════════════════════════════════════════════════════════════════════════════
\appendix
\section{Algorithm details and hyperparameters}\label{app:hyperparams}
%════════════════════════════════════════════════════════════════════════════════

\subsection{Common setup}\label{app:common}

\begin{itemize}
  \item Dual graph $G = (V, E)$ with vertex populations $p_v$ and
    edge attributes (\texttt{shared\_perim} for Iowa).
  \item $k$ districts, $\varepsilon$-population balance per
    Equation~\eqref{eq:pop-balance}, enforced two-sidedly against the
    global ideal on every move (Equation~\eqref{eq:balance-window})
    and audited per step. The bare GerryChain \texttt{recom} proposal
    is wrapped with an explicit whole-plan balance check
    (\texttt{recom\_valid}), since its internal cut search validates
    only one side of each split.
  \item Initial partition from
    \texttt{gerrychain.tree.recursive\_tree\_part}, re-drawn until all
    $k$ districts satisfy the global tolerance
    (\texttt{balanced\_tree\_part}).
  \item Step count: $5 \times 5$: $N = 1\,000$, burn-in $B = 100$;
    Iowa: $N = 500$, $B = 100$; long-run validation
    (\S\ref{sec:exp-longrun}): $2 \times 10^5$ (A) and
    $10^5$ (E) steps.
  \item Random seeds: base seed 42 with a distinct offset per chain
    (independent streams for Python \texttt{random} and
    \texttt{numpy}).
\end{itemize}

\subsection{Chain-specific parameters}\label{app:chains}

\begin{table}[H]
\centering
\small
\begin{tabular}{lll}
\toprule
Chain & Parameter & Value \\
\midrule
A & $\lambda$ (target temperature) & 0.10 \\
B & $p_R$ (ReCom mix probability) & 0.5 \\
B & $\alpha$ (Laplacian weight) & 10.0 \\
B & $\beta$ (fidelity weight) & 1.0 \\
B & Threshold scheme & Cheeger sweep, balance fallback \\
B & Laplacian normalization & False (unnormalized $L$) \\
C & $n_{\mathrm{iqp}}$ (IQP steps per outer) & 10 \\
C & $\alpha$ & 10.0 \\
C & $\beta$ & 1.0 \\
C & MIQP time limit & 5.0 s \\
D & $p_R$ & 0.5 \\
D & Weight randomization & $\mathcal{U}[1, 2]$ per edge \\
D & Repair & False; balanced sweep in place of sign split \\
E & $\alpha$ & 1.0 \\
E & $\beta$ & 10.0 ($\beta \gg \alpha$ keeps $x^\star$ near-binary) \\
E & Rounding & single-shot Bernoulli($x^\star$) \\
E & $\lambda$ (MH target) & 0.10 \\
\bottomrule
\end{tabular}
\caption{Locked-in hyperparameters for the six
  chains.}\label{tab:hyperparams}
\end{table}

\subsection{Diversity controls}\label{app:diversity}

\begin{itemize}
  \item \textbf{Weight jitter}: multiply each edge weight by
    $\mathcal{U}[1, 2]$ before solving (SpecReCom-style). Tested on B
    and C; gave $+27\%$ coverage gain on C with negligible cut quality
    cost. Always on for D (as in Davies et al.'s Algorithm 2). Never
    used in Chain E, whose detailed-balance argument requires the QP
    map to be deterministic given the state and pair.
  \item \textbf{Tabu}: forbid the last $K$ pairs chosen by the local
    kernel (cleared whenever ReCom moves). $K = 3$ in our tests;
    modest cut-tightening, no coverage gain.
\end{itemize}

\subsection{Compute environment}\label{app:compute}

\begin{itemize}
  \item macOS (Darwin 23.6.0), Python 3.12.
  \item Gurobi 11 (academic license) for QP and MIQP.
  \item NetworkX 3.x for graph operations.
  \item Single-threaded; no GPU. All experiments run on a 2024
    MacBook.
\end{itemize}

%════════════════════════════════════════════════════════════════════════════════
\section{Full parameter sweep tables}\label{app:sweeps}
%════════════════════════════════════════════════════════════════════════════════

\needdata{Full sweep tables for $n_{\mathrm{iqp}} \in \{3, 5, 10\}$,
$\alpha \in \{10, 50\}$, $\beta \in \{0, 1\}$, kernel $\in \{k_{\mathrm{unif}},
k_{\mathrm{perim}}, k_{\mathrm{dens}}\}$.}

%════════════════════════════════════════════════════════════════════════════════
\section{Reproducibility checklist}\label{app:reproducibility}
%════════════════════════════════════════════════════════════════════════════════

\begin{itemize}
  \item Code: \texttt{diffusion\_model/recom/} in the project
    repository.
  \item Modules: \texttt{problems.py} (graph setup, balanced initial
    plans), \texttt{chains.py} (chain runners, balance-validated
    ReCom wrapper, per-step audit), \texttt{qp\_model.py}
    (QP/IQP/spectral/QP-MH step factories, global-ideal balance
    windows), \texttt{diagnostics.py} (KL/TV diagnostics).
  \item Orchestrators: \texttt{iowa\_target\_matched.py} (mode and
    seed as CLI arguments) and \texttt{long\_run\_5x5.py}
    (\S\ref{sec:exp-longrun}).
  \item Outputs:
    \texttt{iowa\_hybrid\_test\_out/target\_matched\_*\_seed*.npz},
    each stamped with the full configuration (kernel, $\alpha$,
    $\beta$, $n_{\mathrm{iqp}}$, balance mode, seed) and containing
    per-step cut traces, timings, and balance-audit arrays, so every
    table is regenerable from its own artifact.
  \item Fibers: \texttt{fiber\_5x5/} contains the $5 \times 5$ and
    $4 \times 4$ enumeration scripts and pre-computed
    adjacency files.
  \item Random seeds: base 42, per-chain offsets. Gurobi default seed
    (deterministic given same input).
\end{itemize}

\bibliographystyle{plain}
\bibliography{ref}

\end{document}
